\documentclass[12pt,oneside]{book}

%============================================================
% PACKAGES
%============================================================

%--------------------
% Page + font
%--------------------
\usepackage[margin=1in]{geometry}
\usepackage[T1]{fontenc}
\usepackage[utf8]{inputenc}
\usepackage{lmodern}
\usepackage{microtype}

%--------------------
% Math
%--------------------
\usepackage{amsmath, amssymb, amsfonts, amsthm, mathtools}
\usepackage{bm}
\usepackage{bbm}
\usepackage{mathrsfs}
\usepackage{dsfont}

%--------------------
% Tables + figures
%--------------------
\usepackage{graphicx}
\usepackage{float}
\usepackage{booktabs}
\usepackage{array}
\usepackage{xltabular}
\usepackage{multirow}
\usepackage{caption}
\usepackage{subcaption}

%--------------------
% Lists
%--------------------
\usepackage{enumitem}
\setlist[itemize]{topsep=2pt,itemsep=2pt}
\setlist[enumerate]{topsep=2pt,itemsep=2pt}

%--------------------
% Colors + hyperlinks
%--------------------
\usepackage[dvipsnames]{xcolor}
\usepackage[
    colorlinks=true,
    linkcolor=MidnightBlue,
    citecolor=MidnightBlue,
    urlcolor=MidnightBlue
]{hyperref}
\usepackage[nameinlink,capitalise]{cleveref}

%--------------------
% Algorithms
%--------------------
\usepackage{algorithm}
\usepackage{algpseudocode}

%--------------------
% Code blocks
%--------------------
\usepackage{listings}

\lstdefinestyle{codestyle}{
    basicstyle=\ttfamily\small,
    breaklines=true,
    frame=single,
    numbers=left,
    numberstyle=\tiny,
    keywordstyle=\color{MidnightBlue},
    commentstyle=\color{ForestGreen},
    stringstyle=\color{BrickRed},
    showstringspaces=false,
    tabsize=4
}

\lstset{style=codestyle}

%--------------------
% Section styling
%--------------------
\usepackage{titlesec}
\usepackage{tocloft}

\setcounter{tocdepth}{2}
\setcounter{secnumdepth}{3}

%============================================================
% THEOREMS
%============================================================

\numberwithin{equation}{section}

\theoremstyle{definition}
\newtheorem{definition}{Definition}[section]
\newtheorem{example}{Example}[section]
\newtheorem{exercise}{Exercise}[section]
\newtheorem{remark}{Remark}[section]

\theoremstyle{plain}
\newtheorem{theorem}{Theorem}[section]
\newtheorem{lemma}{Lemma}[section]
\newtheorem{proposition}{Proposition}[section]
\newtheorem{corollary}{Corollary}[section]

\theoremstyle{remark}
\newtheorem*{note}{Note}

%============================================================
% CUSTOM COMMANDS
%============================================================

% Sets
\newcommand{\R}{\mathbb{R}}
\newcommand{\N}{\mathbb{N}}
\newcommand{\Z}{\mathbb{Z}}
\newcommand{\Q}{\mathbb{Q}}
\newcommand{\C}{\mathbb{C}}

% Probability
\newcommand{\E}{\mathbb{E}}
\newcommand{\Var}{\mathrm{Var}}
\newcommand{\Cov}{\mathrm{Cov}}
\newcommand{\Corr}{\mathrm{Corr}}
\newcommand{\Prob}{\mathbb{P}}

\newcommand{\ind}{\mathbbm{1}}
\newcommand{\iid}{\overset{\text{iid}}{\sim}}
\newcommand{\convd}{\overset{d}{\to}}
\newcommand{\convp}{\overset{p}{\to}}
\newcommand{\as}{\overset{a.s.}{\to}}

% Operators
\newcommand{\given}{\,\middle|\,}
\newcommand{\argmin}{\operatorname*{arg\,min}}
\newcommand{\argmax}{\operatorname*{arg\,max}}

% Linear algebra / analysis
\newcommand{\norm}[1]{\left\lVert #1 \right\rVert}
\newcommand{\abs}[1]{\left\lvert #1 \right\rvert}
\newcommand{\inner}[2]{\left\langle #1, #2 \right\rangle}

\newcommand{\dd}{\,\mathrm{d}}
\newcommand{\grad}{\nabla}
\newcommand{\Hess}{\nabla^2}

\newcommand{\tr}{\operatorname{tr}}
\newcommand{\rank}{\operatorname{rank}}
\newcommand{\diag}{\operatorname{diag}}

% Distributions
\newcommand{\Normal}{\mathcal{N}}
\newcommand{\Bern}{\mathrm{Bernoulli}}
\newcommand{\Bin}{\mathrm{Binomial}}
\newcommand{\Pois}{\mathrm{Poisson}}
\newcommand{\Gam}{\mathrm{Gamma}}
\newcommand{\BetaDist}{\mathrm{Beta}}
\newcommand{\Exp}{\mathrm{Exponential}}
\newcommand{\Dir}{\mathrm{Dirichlet}}
\newcommand{\InvGam}{\mathrm{Inv\text{-}Gamma}}

% Title block
\title{EN.553.636 Introduction to Data Science}
\author{Jonathan Y. Ma}
\date{June 2026}

\begin{document}

\maketitle

\noindent \emph{This course was taught by Dr. Tamas Budavari in the Fall 2025 Semester. The notes are transcribed by Jonathan Ma and may contain errors and omitted content, so any issues are to be directed towards me.}

\tableofcontents

\chapter{Descriptive Statistics}

\section{Data Sets}

A data set consisting of $N$ scalar observations is denoted by
\[
\{x_i\}_{i=1}^N, \quad x_i \in \R.
\]

In data analysis, we typically summarize a data set through three main characteristics:
\begin{itemize}
    \item Location (central tendency)
    \item Dispersion (variability)
    \item Shape (distributional features such as skewness and modality)
\end{itemize}

\section{Measures of Location}

\begin{definition}[Mean]
The sample mean of a data set $\{x_i\}_{i=1}^N$ is
\[
\bar{x} = \frac{1}{N} \sum_{i=1}^N x_i.
\]
\end{definition}

In computational settings (e.g., Python), indexing often begins at $0$, in which case the mean is written as
\[
\bar{x} = \frac{1}{N} \sum_{i=0}^{N-1} x_i.
\]

\begin{definition}[Median]
Let $\{x_{(i)}\}_{i=1}^N$ denote the sorted data such that $x_{(1)} \leq \cdots \leq x_{(N)}$.
The sample median is defined as
\[
\text{median} =
\begin{cases}
x_{(\frac{N+1}{2})}, & \text{if } N \text{ is odd}, \\
\frac{1}{2}\big(x_{(\frac{N}{2})} + x_{(\frac{N}{2}+1)}\big), & \text{if } N \text{ is even}.
\end{cases}
\]
\end{definition}

\begin{definition}[Mode]
The mode is the value(s) that occur with the highest frequency in the data set.
A distribution may be unimodal or multimodal depending on the number of such values.
\end{definition}

\section{Measures of Dispersion}

\begin{definition}[Sample Variance]
The sample variance is defined as
\[
s^2 = \frac{1}{N-1} \sum_{i=1}^N (x_i - \bar{x})^2.
\]
\end{definition}

\begin{definition}[Sample Standard Deviation]
The sample standard deviation is
\[
s = \sqrt{s^2}.
\]
\end{definition}

\begin{remark}[Degrees of Freedom]
The factor $N-1$ arises from estimating the population mean using $\bar{x}$.
This reduces the effective degrees of freedom by one, leading to an unbiased estimator of the population variance.
\end{remark}

\begin{proposition}[Unbiasedness of Sample Variance]
Let $X_1, \dots, X_N \iid$ with mean $\mu$ and variance $\sigma^2$. Then
\[
\E[s^2] = \sigma^2.
\]
\end{proposition}

\begin{proof}
We write
\[
\sum_{i=1}^N (X_i - \bar{X})^2
= \sum_{i=1}^N (X_i - \mu)^2 - N(\bar{X} - \mu)^2.
\]
Taking expectations,
\[
\E\left[\sum_{i=1}^N (X_i - \mu)^2\right] = N\sigma^2,
\quad
\E\left[N(\bar{X} - \mu)^2\right] = N \cdot \frac{\sigma^2}{N} = \sigma^2.
\]
Thus,
\[
\E\left[\sum_{i=1}^N (X_i - \bar{X})^2\right] = (N-1)\sigma^2.
\]
Dividing by $N-1$ yields $\E[s^2] = \sigma^2$.
\end{proof}

\section{Outliers}

\begin{definition}[Outlier]
An outlier is an observation that deviates significantly from the rest of the data.
\end{definition}

\begin{remark}
Extreme values can heavily influence the mean and variance. For example, if a single observation diverges (e.g., $x_i \to \infty$), then
\[
\bar{x} \to \infty, \quad s^2 \to \infty.
\]
\end{remark}

\section{Robustness}

\begin{definition}[Robust Statistic]
A statistic is said to be robust if it is not unduly affected by small deviations from model assumptions or the presence of outliers.
\end{definition}

\begin{proposition}
The median is more robust to outliers than the mean.
\end{proposition}

\begin{proof}
The mean depends linearly on all observations:
\[
\bar{x} = \frac{1}{N} \sum_{i=1}^N x_i,
\]
so a single extreme value can arbitrarily shift $\bar{x}$.

In contrast, the median depends only on the ordering of the data. A single extreme value changes only the endpoints of the ordered sequence and does not affect the middle element(s), leaving the median unchanged unless the outlier crosses the median threshold.
\end{proof}

\begin{definition}[Median Absolute Deviation (MAD)]
The median absolute deviation is defined as
\[
\text{MAD} = \mathrm{median} \left( \abs{x_i - \mathrm{median}(x)} \right).
\]
\end{definition}

\begin{remark}
MAD provides a robust alternative to standard deviation, as it is based on medians rather than squared deviations, making it less sensitive to extreme values.
\end{remark}

\chapter{Distributions}

\section{Probability Density Functions}

\begin{definition}[Probability Density Function]
A function $p : \R \to [0,\infty)$ is called a probability density function (PDF) if
\[
\int_{-\infty}^{\infty} p(x)\,dx = 1.
\]
\end{definition}

\begin{proposition}
For a continuous random variable $X$ with density $p(x)$, the probability that $X$ lies in an interval $(a,b)$ is
\[
\Prob(a < X < b) = \int_a^b p(x)\,dx.
\]
\end{proposition}

\begin{example}[Uniform Distribution]
The uniform distribution on $(a,b)$ has density
\[
p(x) = \frac{\ind_{(a,b)}(x)}{b-a},
\]
where $\ind_{(a,b)}(x)$ is the indicator function.
\end{example}

\begin{example}[Gaussian Distribution]
The Gaussian (normal) distribution with mean $\mu$ and variance $\sigma^2$ has density
\[
p(x) = \frac{1}{\sqrt{2\pi\sigma^2}} \exp\left(-\frac{(x-\mu)^2}{2\sigma^2}\right).
\]
\end{example}

\section{Cumulative Distribution Function}

\begin{definition}[Cumulative Distribution Function]
The cumulative distribution function (CDF) of a random variable $X$ is
\[
F(x) = \Prob(X \leq x) = \int_{-\infty}^x p(t)\,dt.
\]
\end{definition}

\begin{proposition}
The CDF $F(x)$ satisfies:
\begin{itemize}
    \item $F(x)$ is non-decreasing,
    \item $\lim_{x \to -\infty} F(x) = 0$ and $\lim_{x \to \infty} F(x) = 1$,
    \item If $p(x)$ is continuous, then $F'(x) = p(x)$.
\end{itemize}
\end{proposition}

\section{Expectation and Moments}

\begin{definition}[Expectation]
The expectation of a random variable $X$ with density $p(x)$ is
\[
\E[X] = \int_{-\infty}^{\infty} x\,p(x)\,dx.
\]
\end{definition}

\begin{definition}[Expectation of a Function]
For any measurable function $f$,
\[
\E[f(X)] = \int_{-\infty}^{\infty} f(x)\,p(x)\,dx.
\]
\end{definition}

\begin{definition}[Moments]
The $k$-th moment of $X$ is
\[
\E[X^k].
\]
The $k$-th central moment is
\[
\E[(X - \mu)^k], \quad \mu = \E[X].
\]
\end{definition}

\begin{definition}[Variance]
The variance of $X$ is
\[
\Var(X) = \E[(X - \mu)^2].
\]
\end{definition}

\begin{proposition}[Variance Identity]
\[
\Var(X) = \E[X^2] - (\E[X])^2.
\]
\end{proposition}

\begin{proof}
Expanding,
\[
\Var(X) = \E[(X - \mu)^2] = \E[X^2 - 2\mu X + \mu^2]
= \E[X^2] - 2\mu \E[X] + \mu^2.
\]
Since $\mu = \E[X]$, this simplifies to
\[
\Var(X) = \E[X^2] - \mu^2.
\]
\end{proof}

\begin{definition}[Standard Deviation]
The standard deviation is
\[
\sigma = \sqrt{\Var(X)}.
\]
\end{definition}

\begin{definition}[Standardized Moments]
The $k$-th standardized moment is
\[
\E\left[\left(\frac{X - \mu}{\sigma}\right)^k\right].
\]
\end{definition}

\begin{definition}[Skewness and Kurtosis]
\begin{itemize}
    \item Skewness: $\E\left[\left(\frac{X - \mu}{\sigma}\right)^3\right]$
    \item Kurtosis: $\E\left[\left(\frac{X - \mu}{\sigma}\right)^4\right]$
\end{itemize}
\end{definition}

\section{Python Introduction}

\subsection{Core Language Concepts}

\begin{itemize}
    \item \textbf{Tuple}: immutable ordered collection
\begin{lstlisting}[language=Python]
x = (1, 2, 3)
\end{lstlisting}

    \item \textbf{List}: mutable ordered collection
\begin{lstlisting}[language=Python]
x = [1, 2, 3]
\end{lstlisting}

    \item \textbf{Function}
\begin{lstlisting}[language=Python]
def f(x):
    return x**2
\end{lstlisting}

    \item \textbf{Lambda}
\begin{lstlisting}[language=Python]
f = lambda x: x**2
\end{lstlisting}

    \item \textbf{For loop}
\begin{lstlisting}[language=Python]
for i in range(5):
    print(i)
\end{lstlisting}

    \item \textbf{Map}
\begin{lstlisting}[language=Python]
list(map(lambda x: x**2, [1,2,3]))
\end{lstlisting}

    \item \textbf{Importing libraries}
\begin{lstlisting}[language=Python]
import numpy as np
import matplotlib.pyplot as plt
\end{lstlisting}
\end{itemize}

\subsection{Pandas Introduction}

Pandas is a library for structured data manipulation.

\begin{lstlisting}[language=Python]
import pandas as pd

df = pd.DataFrame({
    "x": [1,2,3],
    "y": [4,5,6]
})

df.mean()
df.describe()
\end{lstlisting}

\subsection{Seaborn Introduction}

Seaborn is a visualization library built on top of matplotlib.

\begin{lstlisting}[language=Python]
import seaborn as sns

sns.histplot(df["x"])
sns.scatterplot(x="x", y="y", data=df)
\end{lstlisting}

\chapter{Samples and Kernels}

\section{Sample vs Population Quantities}

\begin{proposition}[Sample Estimates vs Population Quantities]
Let $\{x_i\}_{i=1}^N$ be a sample from a random variable $X \sim p(\cdot)$. Then
\[
\bar{x} = \frac{1}{N}\sum_{i=1}^N x_i = \langle x_i \rangle_{i=1}^N
\quad \longleftrightarrow \quad
\mu = \E[X] = \int x\,p(x)\,dx,
\]
\[
s^2 = \frac{1}{N-1}\sum_{i=1}^N (x_i - \bar{x})^2
\quad \longleftrightarrow \quad
\Var(X) = \int (x-\mu)^2 p(x)\,dx.
\]
\end{proposition}

\begin{remark}
Sample statistics approximate population quantities. The quality of approximation improves with increasing $N$.
\end{remark}

\section{Sampling from Distributions}

\begin{proposition}[Uniform Scaling]
If $U \sim \mathrm{Uniform}(0,1)$, then
\[
X = a + (b-a)U \sim \mathrm{Uniform}(a,b).
\]
\end{proposition}

\begin{proof}
For $x \in (a,b)$,
\[
\Prob(X \leq x) = \Prob\left(U \leq \frac{x-a}{b-a}\right) = \frac{x-a}{b-a},
\]
which is the CDF of $\mathrm{Uniform}(a,b)$.
\end{proof}

\begin{proposition}[Inverse Transform Sampling]
Let $U \sim \mathrm{Uniform}(0,1)$ and let $F$ be a strictly increasing CDF. Then
\[
X = F^{-1}(U)
\]
has distribution with CDF $F$.
\end{proposition}

\begin{proof}
For any $x$,
\[
\Prob(X \leq x) = \Prob(F^{-1}(U) \leq x)
= \Prob(U \leq F(x)) = F(x).
\]
\end{proof}

\begin{definition}[Rejection Sampling]
To sample from a target density $p(x)$:
\begin{enumerate}
    \item Choose proposal $q(x)$ and constant $M$ such that $p(x) \leq M q(x)$.
    \item Sample $X \sim q(x)$ and $U \sim \mathrm{Uniform}(0,1)$.
    \item Accept $X$ if $U \leq \frac{p(X)}{M q(X)}$.
\end{enumerate}
\end{definition}

\begin{remark}
Rejection sampling extends naturally to $\R^d$, but efficiency deteriorates in high dimensions due to low acceptance rates.
\end{remark}

\section{Monte Carlo Approximation}

\begin{proposition}[Law of Large Numbers Approximation]
Let $x_i \iid p(\cdot)$. Then
\[
\E[X] = \int x\,p(x)\,dx \approx \frac{1}{N}\sum_{i=1}^N x_i,
\]
\[
\E[f(X)] = \int f(x)\,p(x)\,dx \approx \frac{1}{N}\sum_{i=1}^N f(x_i).
\]
\end{proposition}

\begin{proposition}[Variance Approximation]
\[
\Var(X) = \E[(X-\mu)^2] \approx \frac{1}{N}\sum_{i=1}^N (x_i - \mu)^2.
\]
\end{proposition}

\begin{remark}[Bessel Correction]
The estimator
\[
s^2 = \frac{1}{N-1}\sum_{i=1}^N (x_i - \bar{x})^2
\]
uses $N-1$ degrees of freedom because
\[
\sum_{i=1}^N (x_i - \bar{x}) = 0,
\]
implying only $N-1$ independent deviations.
\end{remark}

\section{Density Estimation}

\subsection{Histograms}

\begin{definition}[Histogram Estimator]
Let bin width be $h$ and bin origin $x_0$. Then
\[
\mathrm{Hist}(x) = \frac{1}{Nh}\sum_{i=1}^N \ind_{\mathrm{bin}(x_i; x_0, h)}(x).
\]
\end{definition}

\begin{remark}
Histograms depend heavily on the choice of bin width $h$ and alignment $x_0$, leading to discontinuous density estimates.
\end{remark}

\subsection{Kernel Density Estimation}

\begin{definition}[Kernel Density Estimator]
Let $K$ be a kernel function and $h>0$ a bandwidth. Then
\[
\widehat{p}(x) = \frac{1}{N}\sum_{i=1}^N K_h(x - x_i)
= \frac{1}{Nh}\sum_{i=1}^N K\left(\frac{x - x_i}{h}\right).
\]
\end{definition}

\begin{remark}
The bandwidth $h$ controls the bias-variance tradeoff:
\begin{itemize}
    \item Small $h$: low bias, high variance (overfitting)
    \item Large $h$: high bias, low variance (oversmoothing)
\end{itemize}
\end{remark}

\begin{definition}[Kernel Function]
A kernel $K$ satisfies:
\begin{itemize}
    \item $K(x) \geq 0$
    \item $\int K(x)\,dx = 1$
    \item Often symmetric: $K(x) = K(-x)$
\end{itemize}
\end{definition}

\begin{example}[Common Kernels]
\begin{itemize}
    \item Uniform: $K(x) = \frac{1}{2}\ind_{\abs{x} \leq 1}$
    \item Triangular: $K(x) = (1-\abs{x})\ind_{\abs{x} \leq 1}$
    \item Gaussian: $K(x) = \frac{1}{\sqrt{2\pi}}e^{-x^2/2}$
    \item Epanechnikov: $K(x) = \frac{3}{4}(1-x^2)\ind_{\abs{x} \leq 1}$
\end{itemize}
\end{example}

\begin{remark}[Support]
Kernels may have:
\begin{itemize}
    \item Finite support (e.g., uniform, Epanechnikov)
    \item Infinite support (e.g., Gaussian)
\end{itemize}
\end{remark}

\begin{remark}[Numerical Considerations]
In practice:
\begin{itemize}
    \item Gaussian kernels are widely used due to smoothness and differentiability
    \item Finite-support kernels are computationally cheaper
    \item Bandwidth selection dominates kernel choice in importance
\end{itemize}
\end{remark}

\chapter{Optimization}

\section{Unconstrained Optimization}

We consider problems of the form
\[
\min_{x \in \R^d} f(x),
\]
where $f$ is typically assumed to be differentiable.

\subsection{First-Order Optimality}

\begin{theorem}[First-Order Condition]
If $f$ is differentiable and $x^\star$ is a local minimizer, then
\[
\grad f(x^\star) = 0.
\]
\end{theorem}

\subsection{Second-Order Conditions}

\begin{theorem}[Second-Order Condition]
If $f$ is twice differentiable and $x^\star$ is a local minimizer, then
\[
\grad f(x^\star) = 0, \quad \Hess f(x^\star) \succeq 0.
\]
If $\Hess f(x^\star) \succ 0$, then $x^\star$ is a strict local minimum.
\end{theorem}

\subsection{Line Search Methods}

Iterative update:
\[
x_{k+1} = x_k + \alpha_k d_k,
\]
where $d_k$ is a descent direction.

\begin{lstlisting}[language=Python]
# Gradient Descent with Line Search
x = x0
for k in range(max_iter):
    grad = grad_f(x)
    d = -grad
    alpha = line_search(f, x, d)
    x = x + alpha * d
\end{lstlisting}

\begin{remark}
Common choices for $d_k$ include $-\grad f(x_k)$ (gradient descent) and quasi-Newton directions.
\end{remark}

\subsection{Newton's Method}

Using second-order information:
\[
x_{k+1} = x_k - \Hess f(x_k)^{-1} \grad f(x_k).
\]

\begin{lstlisting}[language=Python]
# Newton's Method
x = x0
for k in range(max_iter):
    g = grad_f(x)
    H = hess_f(x)
    x = x - np.linalg.solve(H, g)
\end{lstlisting}

\begin{remark}
Newton's method has quadratic convergence near the optimum but requires solving a linear system at each step.
\end{remark}

\subsection{Trust Region Methods}

Solve subproblem:
\[
\min_{\norm{p} \leq \Delta} \quad m_k(p) = f(x_k) + \inner{\grad f(x_k)}{p} + \frac{1}{2} \inner{p}{\Hess f(x_k)p}.
\]

\begin{lstlisting}[language=Python]
# Trust Region (conceptual)
for k in range(max_iter):
    p = solve_trust_region_subproblem(grad, H, Delta)
    if actual_reduction / predicted_reduction > eta:
        x = x + p
        Delta = increase(Delta)
    else:
        Delta = decrease(Delta)
\end{lstlisting}

\section{Constrained Optimization}

General problem:
\[
\min_{x \in \R^d} f(x) \quad \text{s.t.} \quad g_i(x) \leq 0, \; h_j(x) = 0.
\]

\subsection{Karush--Kuhn--Tucker (KKT) Conditions}

\begin{theorem}[KKT Conditions]
If $x^\star$ is optimal and regularity conditions hold, then there exist multipliers $\lambda_i \geq 0$, $\nu_j$ such that
\[
\grad f(x^\star) + \sum_i \lambda_i \grad g_i(x^\star) + \sum_j \nu_j \grad h_j(x^\star) = 0,
\]
\[
g_i(x^\star) \leq 0, \quad h_j(x^\star) = 0,
\]
\[
\lambda_i g_i(x^\star) = 0.
\]
\end{theorem}

\subsection{Simplex Method}

Used for linear programs:
\[
\min c^\top x \quad \text{s.t.} \quad Ax = b, \; x \geq 0.
\]

\begin{lstlisting}[language=Python]
# Conceptual Simplex
initialize basic feasible solution
while optimality not reached:
    choose entering variable
    choose leaving variable
    pivot
\end{lstlisting}

\subsection{Interior Point Methods}

Solve barrier problem:
\[
\min_x \; f(x) - \mu \sum_i \log(-g_i(x)).
\]

\begin{lstlisting}[language=Python]
# Interior Point (Barrier)
x = x0
for t in schedule:
    minimize f(x) - mu * sum(log(-g(x)))
    decrease mu
\end{lstlisting}

\subsection{Ellipsoid Method}

Iteratively shrink ellipsoid containing feasible region.

\begin{lstlisting}[language=Python]
# Ellipsoid (conceptual)
initialize ellipsoid E0
for k in range(max_iter):
    if center feasible:
        update best solution
    else:
        cut ellipsoid using separating hyperplane
\end{lstlisting}

\section{General Optimization Frameworks}

\subsection{Penalty Methods}

Convert constrained problem into unconstrained:
\[
\min_x \; f(x) + \rho \sum_i \max(0, g_i(x))^2.
\]

\begin{remark}
As $\rho \to \infty$, solutions approach feasibility.
\end{remark}

\subsection{Augmented Lagrangian}

\[
\mathcal{L}_\rho(x,\lambda) = f(x) + \lambda^\top g(x) + \frac{\rho}{2}\norm{g(x)}^2.
\]

\begin{lstlisting}[language=Python]
# Augmented Lagrangian
for k in range(max_iter):
    x = minimize L_rho(x, lambda)
    lambda = lambda + rho * g(x)
\end{lstlisting}

\subsection{Alternating Direction Method of Multipliers (ADMM)}

Split problem:
\[
\min f(x) + g(z) \quad \text{s.t.} \quad Ax + Bz = c.
\]

Iterations:
\[
x^{k+1} = \argmin_x \mathcal{L}_\rho(x,z^k,y^k),
\]
\[
z^{k+1} = \argmin_z \mathcal{L}_\rho(x^{k+1},z,y^k),
\]
\[
y^{k+1} = y^k + \rho (Ax^{k+1} + Bz^{k+1} - c).
\]

\begin{lstlisting}[language=Python]
# ADMM
for k in range(max_iter):
    x = argmin_x L(x, z, y)
    z = argmin_z L(x, z, y)
    y = y + rho * (A@x + B@z - c)
\end{lstlisting}

\begin{remark}
ADMM is especially useful for large-scale problems and distributed optimization.
\end{remark}

\chapter{Least Squares}

\section{Multivariate Probability Theory}

\subsection{Dependence}

Let $X,Y \in \R$ be random variables. They are independent if
\[
p(x,y) = p_X(x)p_Y(y).
\]
Equivalently, for events $A$ and $B$,
\[
\Prob(X \in A, Y \in B) = \Prob(X \in A)\Prob(Y \in B).
\]

If
\[
p(x,y) \neq p_X(x)p_Y(y),
\]
then $X$ and $Y$ are dependent.

\subsection{Covariance}

\begin{definition}[Covariance]
The covariance between $X$ and $Y$ is
\[
\Cov(X,Y)
=
\E\left[(X-\E[X])(Y-\E[Y])\right].
\]
\end{definition}

Equivalent notation includes $C_{X,Y}$ and $\sigma(X,Y)$.

\begin{proposition}[Covariance Identity]
\[
\Cov(X,Y) = \E[XY] - \E[X]\E[Y].
\]
\end{proposition}

\begin{proof}
Expanding the definition,
\[
\Cov(X,Y)
=
\E[(X-\mu_X)(Y-\mu_Y)].
\]
Thus,
\[
\Cov(X,Y)
=
\E[XY - X\mu_Y - Y\mu_X + \mu_X\mu_Y].
\]
Using linearity of expectation,
\[
\Cov(X,Y)
=
\E[XY] - \mu_Y\E[X] - \mu_X\E[Y] + \mu_X\mu_Y.
\]
Since $\mu_X=\E[X]$ and $\mu_Y=\E[Y]$,
\[
\Cov(X,Y)
=
\E[XY] - \mu_X\mu_Y.
\]
\end{proof}

\begin{definition}[Sample Covariance]
Given paired data $\{(x_i,y_i)\}_{i=1}^N$, the sample covariance is
\[
C_{xy}
=
\frac{1}{N-1}\sum_{i=1}^N (x_i-\bar{x})(y_i-\bar{y}).
\]
\end{definition}

\subsection{Correlation}

\begin{definition}[Correlation]
The correlation between $X$ and $Y$ is
\[
\Corr(X,Y)
=
\frac{\Cov(X,Y)}{\sigma_X\sigma_Y},
\]
where $\sigma_X=\sqrt{\Var(X)}$ and $\sigma_Y=\sqrt{\Var(Y)}$.
\end{definition}

\begin{remark}
Correlation is a standardized covariance. It satisfies
\[
-1 \leq \Corr(X,Y) \leq 1.
\]
Values near $1$ indicate strong positive linear association, values near $-1$ indicate strong negative linear association, and values near $0$ indicate weak linear association.
\end{remark}

\subsection{Vector Notation}

Let
\[
V=
\begin{pmatrix}
X\\
Y
\end{pmatrix}.
\]
Then
\[
\E[V]
=
\begin{pmatrix}
\E[X]\\
\E[Y]
\end{pmatrix}
=
\begin{pmatrix}
\mu_X\\
\mu_Y
\end{pmatrix}.
\]

\begin{definition}[Covariance Matrix]
The covariance matrix of $V$ is
\[
\Sigma_V
=
\E\left[(V-\E[V])(V-\E[V])^T\right].
\]
For $V=(X,Y)^T$,
\[
\Sigma_V
=
\begin{pmatrix}
\sigma_X^2 & \Cov(X,Y)\\
\Cov(Y,X) & \sigma_Y^2
\end{pmatrix}.
\]
\end{definition}

\begin{remark}
Since $\Cov(X,Y)=\Cov(Y,X)$, covariance matrices are symmetric.
\end{remark}

More generally, for a random vector $X \in \R^d$,
\[
\Sigma
=
\E\left[(X-\mu)(X-\mu)^T\right],
\]
where $\mu=\E[X]$.

\subsection{Bivariate Normal Distribution}

If $X$ and $Y$ are independent normal random variables with means $\mu_X,\mu_Y$ and variances $\sigma_X^2,\sigma_Y^2$, then their joint density is
\[
p(x,y)
=
\frac{1}{2\pi\sigma_X\sigma_Y}
\exp\left[
-\frac{(x-\mu_X)^2}{2\sigma_X^2}
-\frac{(y-\mu_Y)^2}{2\sigma_Y^2}
\right].
\]

In vector notation, for $x \in \R^2$,
\[
p(x)
=
\frac{1}{2\pi |\Sigma|^{1/2}}
\exp\left[
-\frac{1}{2}(x-\mu)^T\Sigma^{-1}(x-\mu)
\right].
\]

Here $|\Sigma|=\det(\Sigma)$.

If
\[
\Sigma
=
\begin{pmatrix}
\sigma_X^2 & 0\\
0 & \sigma_Y^2
\end{pmatrix},
\]
then $X$ and $Y$ are uncorrelated.

\begin{remark}
For jointly Gaussian random variables, uncorrelatedness implies independence. This is not true for arbitrary distributions.
\end{remark}

\subsection{Multivariate Normal Distribution}

\begin{definition}[Multivariate Normal Distribution]
For $x,\mu \in \R^k$ and positive definite covariance matrix $\Sigma \in \R^{k\times k}$,
\[
p(x)
=
\frac{1}{\sqrt{(2\pi)^k|\Sigma|}}
\exp\left[
-\frac{1}{2}(x-\mu)^T\Sigma^{-1}(x-\mu)
\right].
\]
Equivalently,
\[
p(x)
=
\frac{1}{\sqrt{|2\pi\Sigma|}}
\exp\left[
-\frac{1}{2}(x-\mu)^T\Sigma^{-1}(x-\mu)
\right].
\]
\end{definition}

\section{Ordinary Least Squares}

\subsection{The Idea}

Suppose we have a training set
\[
\{(x_i,y_i)\}_{i=1}^N.
\]
We want to fit a parameterized function
\[
f(x;\theta),
\]
where $\theta$ may contain one or more parameters.

The residual for observation $i$ is
\[
r_i(\theta)=y_i-f(x_i;\theta).
\]

Ordinary least squares chooses parameters by minimizing the sum of squared residuals:
\[
\hat{\theta}
=
\argmin_{\theta}
\sum_{i=1}^N
\left[y_i-f(x_i;\theta)\right]^2.
\]

\subsection{Fitting a Constant}

Consider the constant model
\[
f(x;\mu)=\mu.
\]
The cost function is
\[
C(\mu)
=
\sum_{i=1}^N (y_i-\mu)^2.
\]

\begin{proposition}[Least Squares Estimate of a Constant]
The least squares estimate of $\mu$ is the sample mean:
\[
\hat{\mu}
=
\frac{1}{N}\sum_{i=1}^N y_i.
\]
\end{proposition}

\begin{proof}
Differentiate the cost:
\[
\frac{dC}{d\mu}
=
\frac{d}{d\mu}
\sum_{i=1}^N (y_i-\mu)^2
=
-2\sum_{i=1}^N (y_i-\mu).
\]
At the optimum,
\[
-2\sum_{i=1}^N (y_i-\hat{\mu})=0.
\]
Therefore,
\[
\sum_{i=1}^N y_i - N\hat{\mu}=0,
\]
and hence
\[
\hat{\mu}
=
\frac{1}{N}\sum_{i=1}^N y_i.
\]
\end{proof}

\subsection{Weighted Constant Fit and Heteroscedasticity}

Suppose observations have different variances $\sigma_i^2$. Define weights
\[
w_i=\frac{1}{\sigma_i^2}.
\]
The weighted least squares cost is
\[
C(\mu)
=
\sum_{i=1}^N w_i(y_i-\mu)^2.
\]

\begin{proposition}[Weighted Average]
The weighted least squares estimate is
\[
\hat{\mu}
=
\frac{\sum_{i=1}^N w_i y_i}{\sum_{i=1}^N w_i}.
\]
\end{proposition}

\begin{proof}
Differentiate:
\[
\frac{dC}{d\mu}
=
-2\sum_{i=1}^N w_i(y_i-\mu).
\]
At the optimum,
\[
\sum_{i=1}^N w_i(y_i-\hat{\mu})=0.
\]
Thus,
\[
\sum_{i=1}^N w_i y_i
-
\hat{\mu}\sum_{i=1}^N w_i
=
0.
\]
Solving for $\hat{\mu}$ gives
\[
\hat{\mu}
=
\frac{\sum_{i=1}^N w_i y_i}{\sum_{i=1}^N w_i}.
\]
\end{proof}

\section{Simple Linear Regression}

Consider the linear model
\[
f(x;\theta)=a+bx,
\]
where
\[
\theta=
\begin{pmatrix}
a\\
b
\end{pmatrix}.
\]
The least squares estimate solves
\[
\hat{\theta}
=
\argmin_{a,b}
\sum_{i=1}^N
\left[y_i-(a+bx_i)\right]^2.
\]

The cost is quadratic in $a$ and $b$, so the first-order conditions form a linear system.

\begin{proposition}[Simple Linear Regression Coefficients]
The least squares estimates are
\[
\hat{b}
=
\frac{\sum_{i=1}^N (x_i-\bar{x})(y_i-\bar{y})}
{\sum_{i=1}^N (x_i-\bar{x})^2},
\]
and
\[
\hat{a}
=
\bar{y}-\hat{b}\bar{x}.
\]
\end{proposition}

\begin{proof}
The cost is
\[
C(a,b)
=
\sum_{i=1}^N
(y_i-a-bx_i)^2.
\]
The first-order conditions are
\[
\frac{\partial C}{\partial a}
=
-2\sum_{i=1}^N (y_i-a-bx_i)
=
0,
\]
and
\[
\frac{\partial C}{\partial b}
=
-2\sum_{i=1}^N x_i(y_i-a-bx_i)
=
0.
\]
The first equation gives
\[
\sum_{i=1}^N y_i
=
Na+b\sum_{i=1}^N x_i.
\]
Dividing by $N$,
\[
\bar{y}=a+b\bar{x},
\]
so
\[
a=\bar{y}-b\bar{x}.
\]
Substituting into the second first-order condition yields
\[
\sum_{i=1}^N x_i\left[y_i-\bar{y}-b(x_i-\bar{x})\right]=0.
\]
Equivalently,
\[
\sum_{i=1}^N (x_i-\bar{x})(y_i-\bar{y})
=
b\sum_{i=1}^N (x_i-\bar{x})^2.
\]
Thus,
\[
\hat{b}
=
\frac{\sum_{i=1}^N (x_i-\bar{x})(y_i-\bar{y})}
{\sum_{i=1}^N (x_i-\bar{x})^2},
\]
and
\[
\hat{a}=\bar{y}-\hat{b}\bar{x}.
\]
\end{proof}

\section{Linear Regression with Basis Functions}

Let $\phi_1,\dots,\phi_K$ be known basis functions. A linear regression model is
\[
f(x;\beta)
=
\sum_{k=1}^K \beta_k \phi_k(x).
\]
Equivalently,
\[
f(x;\beta)=\beta^T\phi(x),
\]
where
\[
\beta=(\beta_1,\dots,\beta_K)^T,
\quad
\phi(x)=(\phi_1(x),\dots,\phi_K(x))^T.
\]

The cost function is
\[
C(\beta)
=
\sum_{i=1}^N
\left[
y_i-\sum_{k=1}^K \beta_k\phi_k(x_i)
\right]^2.
\]

Define the design matrix $X \in \R^{N\times K}$ by
\[
X_{ik}=\phi_k(x_i).
\]
Then the model predictions are
\[
\hat{y}=X\beta,
\]
and the cost becomes
\[
C(\beta)
=
\norm{y-X\beta}_2^2.
\]

\subsection{Normal Equations}

\begin{theorem}[Normal Equations]
The least squares solution satisfies
\[
X^T X\hat{\beta}=X^T y.
\]
If $X^T X$ is invertible, then
\[
\hat{\beta}
=
(X^T X)^{-1}X^T y.
\]
\end{theorem}

\begin{proof}
The cost is
\[
C(\beta)
=
(y-X\beta)^T(y-X\beta).
\]
Expanding,
\[
C(\beta)
=
y^Ty - 2\beta^T X^T y + \beta^T X^T X\beta.
\]
Taking the gradient with respect to $\beta$,
\[
\grad_\beta C(\beta)
=
-2X^Ty + 2X^T X\beta.
\]
At the optimum,
\[
-2X^Ty + 2X^T X\hat{\beta}=0.
\]
Therefore,
\[
X^T X\hat{\beta}=X^T y.
\]
If $X^T X$ is invertible,
\[
\hat{\beta}
=
(X^T X)^{-1}X^T y.
\]
\end{proof}

\begin{remark}[Kronecker Delta]
The Kronecker delta is
\[
\delta_{kl}
=
\begin{cases}
1, & k=l,\\
0, & k\neq l.
\end{cases}
\]
It satisfies
\[
\sum_l \delta_{kl}a_l=a_k.
\]
This is the componentwise analogue of
\[
Ia=a.
\]
\end{remark}

\subsection{Componentwise Derivation}

Starting from
\[
C(\beta)
=
\sum_{i=1}^N
\left[
y_i-\sum_{k=1}^K X_{ik}\beta_k
\right]^2,
\]
differentiate with respect to $\beta_l$:
\[
\frac{\partial C}{\partial \beta_l}
=
-2\sum_{i=1}^N
\left[
y_i-\sum_{k=1}^K X_{ik}\beta_k
\right]X_{il}.
\]
At the optimum,
\[
\sum_{i=1}^N
\left[
y_i-\sum_{k=1}^K X_{ik}\hat{\beta}_k
\right]X_{il}
=
0.
\]
Therefore,
\[
\sum_{i=1}^N X_{il}y_i
=
\sum_{k=1}^K
\left(
\sum_{i=1}^N X_{il}X_{ik}
\right)\hat{\beta}_k.
\]
This is precisely
\[
X^Ty=X^TX\hat{\beta}.
\]

\subsection{Moore--Penrose Pseudoinverse}

When $X^TX$ is invertible,
\[
X^+ = (X^TX)^{-1}X^T,
\]
so
\[
\hat{\beta}=X^+y.
\]

More generally, the Moore--Penrose pseudoinverse can be defined using the singular value decomposition. If
\[
X=U\Sigma V^T,
\]
then
\[
X^+=V\Sigma^+U^T,
\]
where $\Sigma^+$ is formed by taking reciprocals of the nonzero singular values.

\section{Geometry of Least Squares}

Least squares projects $y$ onto the column space of $X$.

The fitted values are
\[
\hat{y}=X\hat{\beta}.
\]
Using the closed form solution,
\[
\hat{y}
=
X(X^TX)^{-1}X^Ty.
\]

\begin{definition}[Hat Matrix]
The hat matrix is
\[
H
=
X(X^TX)^{-1}X^T.
\]
Thus,
\[
\hat{y}=Hy.
\]
\end{definition}

\begin{proposition}[Projection Properties]
The hat matrix satisfies
\[
H^T=H,
\]
and
\[
H^2=H.
\]
\end{proposition}

\begin{proof}
First,
\[
H^T
=
\left[X(X^TX)^{-1}X^T\right]^T
=
X(X^TX)^{-1}X^T
=
H,
\]
since $X^TX$ is symmetric.

Next,
\[
H^2
=
X(X^TX)^{-1}X^T X(X^TX)^{-1}X^T.
\]
Since $X^TX$ appears in the middle,
\[
H^2
=
X(X^TX)^{-1}(X^TX)(X^TX)^{-1}X^T
=
X(X^TX)^{-1}X^T
=
H.
\]
\end{proof}

\begin{remark}
The residual vector is
\[
r=y-\hat{y}=y-X\hat{\beta}.
\]
The normal equations imply
\[
X^Tr=0.
\]
Thus, the residuals are orthogonal to every column of the design matrix.
\end{remark}

\section{Weighted Least Squares}

For heteroscedastic errors, use
\[
C(\beta)
=
\sum_{i=1}^N
w_i
\left[
y_i-\sum_{k=1}^K X_{ik}\beta_k
\right]^2.
\]
Let $W$ be diagonal with
\[
W_{ij}=w_i\delta_{ij}.
\]
Then
\[
C(\beta)
=
(y-X\beta)^TW(y-X\beta).
\]

\begin{theorem}[Weighted Least Squares]
The weighted least squares estimator satisfies
\[
X^TWX\hat{\beta}=X^TWy.
\]
If $X^TWX$ is invertible, then
\[
\hat{\beta}
=
(X^TWX)^{-1}X^TWy.
\]
\end{theorem}

\begin{proof}
Expand the weighted cost:
\[
C(\beta)
=
(y-X\beta)^TW(y-X\beta).
\]
Taking the gradient,
\[
\grad_\beta C(\beta)
=
-2X^TWy+2X^TWX\beta.
\]
Setting the gradient equal to zero gives
\[
X^TWX\hat{\beta}=X^TWy.
\]
Solving gives
\[
\hat{\beta}
=
(X^TWX)^{-1}X^TWy.
\]
\end{proof}

\begin{remark}
If errors are uncorrelated but have variances $\sigma_i^2$, then the optimal weights are
\[
w_i=\frac{1}{\sigma_i^2}.
\]
If errors are correlated with covariance matrix $C$, then the natural weight matrix is
\[
W=C^{-1}.
\]
\end{remark}

\section{Regularization}

Least squares can become unstable when predictors are highly correlated or when $X^TX$ is nearly singular. Regularization stabilizes estimation by penalizing large coefficients.

\subsection{Ridge Regression}

\begin{definition}[Ridge Regression]
Ridge regression solves
\[
\hat{\beta}
=
\argmin_{\beta}
\norm{y-X\beta}_2^2
+
\lambda\norm{\beta}_2^2,
\]
where $\lambda \geq 0$.
\end{definition}

\begin{theorem}[Ridge Solution]
The ridge regression estimator is
\[
\hat{\beta}_{\mathrm{ridge}}
=
(X^TX+\lambda I)^{-1}X^Ty.
\]
\end{theorem}

\begin{proof}
The objective is
\[
C(\beta)
=
(y-X\beta)^T(y-X\beta)+\lambda\beta^T\beta.
\]
Taking the gradient,
\[
\grad_\beta C(\beta)
=
-2X^Ty+2X^TX\beta+2\lambda\beta.
\]
Setting the gradient equal to zero,
\[
(X^TX+\lambda I)\hat{\beta}=X^Ty.
\]
Thus,
\[
\hat{\beta}_{\mathrm{ridge}}
=
(X^TX+\lambda I)^{-1}X^Ty.
\]
\end{proof}

More generally, one may use a matrix penalty:
\[
\hat{\beta}
=
\argmin_{\beta}
\norm{y-X\beta}_2^2
+
\lambda\norm{\Gamma\beta}_2^2.
\]
The solution satisfies
\[
(X^TX+\lambda\Gamma^T\Gamma)\hat{\beta}=X^Ty.
\]

\subsection{Lasso Regression}

\begin{definition}[Lasso Regression]
Lasso regression solves
\[
\hat{\beta}
=
\argmin_{\beta}
\norm{y-X\beta}_2^2
+
\lambda\norm{\beta}_1.
\]
\end{definition}

\begin{remark}
The $L_1$ penalty encourages sparsity, meaning some coefficients are exactly zero. Unlike ridge regression, lasso does not generally have a simple closed form solution.
\end{remark}

\section{Least Squares Algorithm}

\begin{algorithm}[H]
\caption{Ordinary Least Squares via Normal Equations}
\begin{algorithmic}[1]
\Require Design matrix $X \in \R^{N\times K}$, response vector $y \in \R^N$
\Ensure Coefficient estimate $\hat{\beta}$
\State Compute $A \gets X^TX$
\State Compute $b \gets X^Ty$
\State Solve $A\hat{\beta}=b$
\State \Return $\hat{\beta}$
\end{algorithmic}
\end{algorithm}

\begin{algorithm}[H]
\caption{Weighted Least Squares}
\begin{algorithmic}[1]
\Require Design matrix $X$, response vector $y$, weights $w_1,\dots,w_N$
\Ensure Weighted least squares estimate $\hat{\beta}$
\State Form diagonal matrix $W$ with $W_{ii}=w_i$
\State Compute $A \gets X^TWX$
\State Compute $b \gets X^TWy$
\State Solve $A\hat{\beta}=b$
\State \Return $\hat{\beta}$
\end{algorithmic}
\end{algorithm}

\section{Implementation Notes}

\begin{lstlisting}[language=Python]
import numpy as np

def ordinary_least_squares(X, y):
    """
    Computes beta_hat = (X^T X)^(-1) X^T y.
    Uses solve rather than explicitly forming the inverse.
    """
    A = X.T @ X
    b = X.T @ y
    beta_hat = np.linalg.solve(A, b)
    return beta_hat
\end{lstlisting}

\begin{lstlisting}[language=Python]
def weighted_least_squares(X, y, w):
    """
    Weighted least squares with diagonal weights w.
    """
    W = np.diag(w)
    A = X.T @ W @ X
    b = X.T @ W @ y
    beta_hat = np.linalg.solve(A, b)
    return beta_hat
\end{lstlisting}

\begin{lstlisting}[language=Python]
def ridge_regression(X, y, lam):
    """
    Ridge regression:
    beta_hat = (X^T X + lambda I)^(-1) X^T y.
    """
    K = X.shape[1]
    A = X.T @ X + lam * np.eye(K)
    b = X.T @ y
    beta_hat = np.linalg.solve(A, b)
    return beta_hat
\end{lstlisting}

\section{Linear Combinations and Basis Expansions}

Linear regression is fundamentally about representing data as linear combinations of known basis functions:
\[
f(x;\beta)
=
\sum_{k=1}^K \beta_k\phi_k(x).
\]

Examples include:
\begin{itemize}
    \item Polynomial regression
    \item Fourier series
    \item Discrete cosine transforms
    \item Spherical harmonics
    \item Image bases
    \item Shape bases
\end{itemize}

\begin{remark}
The term ``linear regression'' means linear in the parameters $\beta$, not necessarily linear in the input $x$. For example,
\[
f(x;\beta)=\beta_0+\beta_1x+\beta_2x^2
\]
is still a linear regression model because it is linear in $\beta$.
\end{remark}

\chapter{Principal Components Analysis}

\section{Directions of Maximum Variance}

Let $X \in \R^N$ be a random vector with
\[
\E[X]=0
\]
and covariance matrix
\[
C=\E[XX^T].
\]
For any direction $a \in \R^N$, the scalar projection of $X$ onto $a$ is
\[
a^TX.
\]
The variance of this projection is
\[
\Var(a^TX)
=
\E[(a^TX)^2].
\]
Since
\[
(a^TX)^2 = a^TXX^Ta,
\]
we have
\[
\Var(a^TX)
=
\E[a^TXX^Ta]
=
a^T\E[XX^T]a
=
a^TCa.
\]

Therefore, the direction of maximum variance solves
\[
\max_{a \in \R^N} a^TCa
\quad
\text{s.t.}
\quad
a^Ta=1.
\]

\section{Constrained Optimization Derivation}

Define the Lagrangian
\[
\mathcal{L}(a,\lambda)
=
a^TCa-\lambda(a^Ta-1).
\]
The constrained maximizer satisfies
\[
\frac{\partial \mathcal{L}}{\partial \lambda}=0,
\quad
\grad_a \mathcal{L}=0.
\]
The first condition gives
\[
a^Ta=1.
\]

For the second condition, write
\[
a^TCa
=
\sum_{i=1}^N\sum_{j=1}^N a_i C_{ij}a_j.
\]
Differentiating with respect to $a_k$,
\[
\frac{\partial}{\partial a_k}
\left(
\sum_{i=1}^N\sum_{j=1}^N a_iC_{ij}a_j
\right)
=
\sum_{i=1}^N\sum_{j=1}^N
\frac{\partial a_i}{\partial a_k}C_{ij}a_j
+
\sum_{i=1}^N\sum_{j=1}^N
a_iC_{ij}\frac{\partial a_j}{\partial a_k}.
\]
Using
\[
\frac{\partial a_i}{\partial a_k}=\delta_{ik},
\]
we get
\[
\sum_{j=1}^N C_{kj}a_j
+
\sum_{i=1}^N a_iC_{ik}.
\]
In matrix notation, this is
\[
Ca+C^Ta.
\]
Since covariance matrices are symmetric, $C=C^T$, so
\[
\grad_a(a^TCa)=2Ca.
\]
Also,
\[
\grad_a\{\lambda(a^Ta-1)\}=2\lambda a.
\]
Thus,
\[
2Ca-2\lambda a=0,
\]
so
\[
Ca=\lambda a.
\]

\begin{theorem}[First Principal Component]
The direction of maximum variance is the eigenvector $a_1$ corresponding to the largest eigenvalue $\lambda_1$ of $C$:
\[
Ca_1=\lambda_1a_1.
\]
The maximum variance is
\[
a_1^TCa_1=\lambda_1.
\]
\end{theorem}

\begin{proof}
From the Lagrange multiplier derivation, any stationary direction satisfies
\[
Ca=\lambda a.
\]
If $a$ is normalized so that $a^Ta=1$, then
\[
a^TCa
=
a^T\lambda a
=
\lambda a^Ta
=
\lambda.
\]
Therefore, maximizing $a^TCa$ over unit vectors is equivalent to choosing the largest eigenvalue. The associated eigenvector is the first principal component.
\end{proof}

\section{Second Principal Component}

The second principal component is the direction of maximum variance subject to being uncorrelated with the first principal component.

Let $a_1$ be the first principal component. Since
\[
\Cov(a^TX,a_1^TX)
=
a^TCa_1,
\]
the uncorrelatedness constraint is
\[
a^TCa_1=0.
\]

The constrained optimization problem is
\[
\max_{a \in \R^N}
\left[
a^TCa
-
\lambda(a^Ta-1)
-
\lambda'(a^TCa_1)
\right].
\]

The first-order condition is
\[
2Ca-2\lambda a-\lambda'Ca_1=0.
\]
Since
\[
Ca_1=\lambda_1a_1,
\]
left multiply by $a_1^T$:
\[
2a_1^TCa
-
2\lambda a_1^Ta
-
\lambda'a_1^TCa_1
=
0.
\]
The orthogonality and uncorrelatedness constraints imply
\[
a_1^TCa=0,
\quad
a_1^Ta=0.
\]
Therefore,
\[
-\lambda'a_1^TCa_1=0.
\]
But
\[
a_1^TCa_1=\lambda_1,
\]
so
\[
\lambda'=0.
\]
Thus the first-order condition reduces to
\[
Ca=\lambda a.
\]

\begin{theorem}[Second Principal Component]
The second principal component is the eigenvector $a_2$ corresponding to the second largest eigenvalue $\lambda_2$:
\[
Ca_2=\lambda_2a_2.
\]
\end{theorem}

\section{Principal Components Analysis}

Let
\[
\lambda_1 \geq \lambda_2 \geq \cdots \geq \lambda_N \geq 0
\]
be the eigenvalues of $C$, with corresponding orthonormal eigenvectors
\[
e_1,\dots,e_N.
\]
Since $C$ is symmetric positive semidefinite, it admits the spectral decomposition
\[
C
=
\sum_{k=1}^N \lambda_k e_ke_k^T.
\]

For any eigenvector $e_l$,
\[
Ce_l
=
\sum_{k=1}^N \lambda_k e_k(e_k^Te_l).
\]
Since the eigenvectors are orthonormal,
\[
e_k^Te_l=\delta_{kl}.
\]
Therefore,
\[
Ce_l
=
\lambda_l e_l.
\]

In matrix form, let
\[
E=
\begin{pmatrix}
e_1 & e_2 & \cdots & e_N
\end{pmatrix},
\]
and
\[
\Lambda=
\diag(\lambda_1,\dots,\lambda_N).
\]
Then
\[
C=E\Lambda E^T.
\]

\begin{definition}[Principal Components]
The principal components are the orthonormal eigenvectors of the covariance matrix, ordered by decreasing eigenvalue.
\end{definition}

\begin{remark}
The eigenvectors associated with the largest eigenvalues capture the directions of greatest variance in the data.
\end{remark}

\section{Low-Rank Approximation}

If only $K<N$ principal components are retained, then
\[
C
\approx
C_K
=
\sum_{k=1}^K \lambda_k e_ke_k^T.
\]
Equivalently,
\[
C_K
=
E_K\Lambda_K E_K^T,
\]
where
\[
E_K=
\begin{pmatrix}
e_1 & \cdots & e_K
\end{pmatrix}
\]
and
\[
\Lambda_K=\diag(\lambda_1,\dots,\lambda_K).
\]

The remaining covariance is
\[
C-C_K
=
\sum_{l=K+1}^N \lambda_l e_le_l^T.
\]

\begin{definition}[Projection Matrix]
The rank-$K$ PCA projection matrix is
\[
P_K=E_KE_K^T.
\]
\end{definition}

\begin{remark}
The matrix $P_K$ projects observations onto the subspace spanned by the first $K$ principal components.
\end{remark}

\section{New Coordinate System}

Since $E$ is orthonormal,
\[
E^TE=EE^T=I.
\]
Thus, the transformation
\[
Z=E^TX
\]
expresses $X$ in the principal component coordinate system.

Using only the first $K$ eigenvectors,
\[
Z_K=E_K^TX.
\]
The reconstruction in the original coordinates is
\[
X_K=E_KZ_K.
\]
Therefore,
\[
X_K=E_KE_K^TX=P_KX.
\]

\begin{remark}
PCA can be viewed as a rotation followed by truncation. The full transformation $E^TX$ preserves all information, while the truncated transformation $E_K^TX$ discards directions with smaller variance.
\end{remark}

\section{Samples and the Sample Covariance Matrix}

Suppose we have $n$ centered observations
\[
x_1,\dots,x_n \in \R^N,
\]
arranged as columns of a data matrix
\[
X=
\begin{pmatrix}
x_1 & x_2 & \cdots & x_n
\end{pmatrix}.
\]
Assuming the sample mean is zero, the sample covariance matrix is
\[
C
=
\frac{1}{n-1}\sum_{i=1}^n x_ix_i^T.
\]
In matrix form,
\[
C
=
\frac{1}{n-1}XX^T.
\]

\section{Connection to Singular Value Decomposition}

Let the singular value decomposition of $X$ be
\[
X=UWV^T,
\]
where
\[
U^TU=I,
\quad
V^TV=I,
\]
and $W$ is diagonal with nonnegative singular values.

Then
\[
C
=
\frac{1}{n-1}XX^T.
\]
Substituting the SVD,
\[
C
=
\frac{1}{n-1}
UWV^TVWU^T.
\]
Since
\[
V^TV=I,
\]
we obtain
\[
C
=
\frac{1}{n-1}UW^2U^T.
\]

Thus, if
\[
C=E\Lambda E^T,
\]
then
\[
E=U,
\]
and
\[
\Lambda=\frac{1}{n-1}W^2.
\]

\begin{remark}
Computing PCA by SVD is often more numerically stable than explicitly forming the covariance matrix.
\end{remark}

\section{Whitening}

\begin{definition}[Whitening Transformation]
Given covariance decomposition
\[
C=E\Lambda E^T,
\]
the whitened representation of centered data $X$ is
\[
A=\Lambda^{-1/2}E^TX.
\]
\end{definition}

\begin{proposition}
If $X$ has covariance $C=E\Lambda E^T$, then the whitened variable
\[
A=\Lambda^{-1/2}E^TX
\]
has identity covariance.
\end{proposition}

\begin{proof}
Compute
\[
\Cov(A)
=
\E[AA^T].
\]
Substituting $A=\Lambda^{-1/2}E^TX$,
\[
\Cov(A)
=
\E[\Lambda^{-1/2}E^TXX^TE\Lambda^{-1/2}].
\]
Since $\E[XX^T]=C$,
\[
\Cov(A)
=
\Lambda^{-1/2}E^TCE\Lambda^{-1/2}.
\]
Using $C=E\Lambda E^T$,
\[
\Cov(A)
=
\Lambda^{-1/2}E^TE\Lambda E^TE\Lambda^{-1/2}.
\]
Since $E^TE=I$,
\[
\Cov(A)
=
\Lambda^{-1/2}\Lambda\Lambda^{-1/2}
=
I.
\]
\end{proof}

\section{Scree Plot and Explained Variance}

The eigenvalue spectrum is
\[
\{\lambda_1,\lambda_2,\dots,\lambda_N\}.
\]
A scree plot displays these eigenvalues in decreasing order.

\begin{definition}[Explained Variance Ratio]
The explained variance ratio of component $k$ is
\[
\frac{\lambda_k}{\sum_{j=1}^N \lambda_j}.
\]
The cumulative explained variance of the first $K$ components is
\[
\frac{\sum_{k=1}^K \lambda_k}{\sum_{j=1}^N \lambda_j}.
\]
\end{definition}

\begin{remark}
A common heuristic is to choose $K$ by looking for an elbow in the scree plot. The elbow indicates that additional components contribute relatively little variance.
\end{remark}

\begin{proposition}[Variance Additivity for Uncorrelated Components]
If $Z_1,\dots,Z_K$ are pairwise uncorrelated, then
\[
\Var\left(\sum_{k=1}^K Z_k\right)
=
\sum_{k=1}^K \Var(Z_k).
\]
\end{proposition}

\begin{proof}
For two variables,
\[
\Var(X+Y)
=
\Var(X)+\Var(Y)+2\Cov(X,Y).
\]
If $X$ and $Y$ are uncorrelated, then $\Cov(X,Y)=0$, so
\[
\Var(X+Y)=\Var(X)+\Var(Y).
\]
The general result follows by applying the same expansion to $K$ variables:
\[
\Var\left(\sum_{k=1}^K Z_k\right)
=
\sum_{k=1}^K \Var(Z_k)
+
2\sum_{i<j}\Cov(Z_i,Z_j).
\]
If all variables are pairwise uncorrelated, all covariance terms vanish.
\end{proof}

\section{PCA Algorithm}

\begin{algorithm}[H]
\caption{Principal Components Analysis}
\begin{algorithmic}[1]
\Require Data matrix $X \in \R^{n\times d}$ with rows as observations, number of components $K$
\Ensure Principal component scores $Z_K$, loading matrix $E_K$
\State Center the data: $X_c \gets X - \bar{X}$
\State Compute covariance matrix: $C \gets \frac{1}{n-1}X_c^TX_c$
\State Compute eigendecomposition: $C=E\Lambda E^T$
\State Sort eigenvalues in decreasing order and reorder eigenvectors accordingly
\State Keep first $K$ eigenvectors: $E_K \gets [e_1,\dots,e_K]$
\State Project data: $Z_K \gets X_cE_K$
\State \Return $Z_K, E_K$
\end{algorithmic}
\end{algorithm}

\begin{algorithm}[H]
\caption{PCA via Singular Value Decomposition}
\begin{algorithmic}[1]
\Require Centered data matrix $X_c \in \R^{n\times d}$, number of components $K$
\Ensure Principal component scores $Z_K$, loading matrix $V_K$
\State Compute SVD: $X_c=U W V^T$
\State Keep first $K$ right singular vectors: $V_K \gets [v_1,\dots,v_K]$
\State Compute scores: $Z_K \gets X_cV_K$
\State Eigenvalues are $\lambda_k = \frac{w_k^2}{n-1}$
\State \Return $Z_K,V_K$
\end{algorithmic}
\end{algorithm}

\section{Implementation Notes}

\begin{lstlisting}[language=Python]
import numpy as np

def pca(X, K):
    """
    PCA using SVD.
    Rows are observations and columns are features.
    """
    X_centered = X - X.mean(axis=0)

    U, S, Vt = np.linalg.svd(X_centered, full_matrices=False)

    components = Vt[:K]
    scores = X_centered @ components.T
    eigenvalues = (S**2) / (X.shape[0] - 1)

    explained_variance_ratio = eigenvalues / eigenvalues.sum()

    return scores, components, eigenvalues[:K], explained_variance_ratio[:K]
\end{lstlisting}

\begin{remark}
For machine learning data matrices, the common convention is rows as observations and columns as features. In this convention, the covariance matrix is
\[
C=\frac{1}{n-1}X_c^TX_c.
\]
In the column-vector convention used in the derivation above, the covariance matrix is
\[
C=\frac{1}{n-1}X_cX_c^T.
\]
The mathematics is the same, but the orientation of the data matrix changes.
\end{remark}

\chapter{Nearest Neighbors}

\section{Classification}

Classification is the supervised learning problem of assigning discrete labels to observations.

A training set consists of labeled data:
\[
T=\{(x_i,C_i)\}_{i=1}^N,
\]
where
\[
x_i \in \R^d
\]
is a feature vector and $C_i$ is the known class label.

A query set consists of unlabeled data:
\[
Q=\{x_i\}_{i=1}^M,
\]
where
\[
x_i \in \R^d.
\]

The goal is to construct a classifier
\[
\hat{C}(x)
\]
that assigns a class label to a new point $x$.

\begin{remark}
Classification is similar to regression, but the response variable is discrete rather than continuous. Common classifiers include nearest neighbors, naive Bayes, LDA/QDA, logistic regression, decision trees, random forests, and support vector machines.
\end{remark}

\section{Nearest Neighbors}

\begin{definition}[Nearest Neighbor Classifier]
Given a query point $x$, the nearest neighbor classifier assigns $x$ the label of the closest training point:
\[
\hat{C}(x)=C_{i^\star},
\]
where
\[
i^\star=\argmin_{1\leq i\leq N} d(x,x_i).
\]
\end{definition}

The most common distance is Euclidean distance:
\[
d(x,x_i)=\norm{x-x_i}_2
=
\sqrt{\sum_{j=1}^d (x_j-x_{ij})^2}.
\]

\section{$k$-Nearest Neighbors}

\begin{definition}[$k$-Nearest Neighbor Classifier]
Let $\mathcal{N}_k(x)$ denote the indices of the $k$ nearest training points to $x$. The $k$-nearest neighbor classifier assigns the most frequent class among those neighbors:
\[
\hat{C}(x)
=
\argmax_{c}
\sum_{i\in \mathcal{N}_k(x)}
\ind(C_i=c).
\]
\end{definition}

A weighted version uses distance-dependent weights:
\[
\hat{C}(x)
=
\argmax_c
\sum_{i\in \mathcal{N}_k(x)}
w_i(x)\ind(C_i=c),
\]
where a common choice is
\[
w_i(x)=\frac{1}{d(x,x_i)+\varepsilon}.
\]

\begin{remark}
Using $k$ instead of a fixed distance cutoff adapts better to regions with different data density. In dense regions, the neighborhood is small; in sparse regions, the neighborhood expands.
\end{remark}

\section{Bias-Variance Tradeoff}

The choice of $k$ controls model complexity.

\begin{itemize}
    \item Small $k$: low bias, high variance
    \item Large $k$: high bias, low variance
\end{itemize}

\begin{remark}
The case $k=1$ produces highly flexible decision boundaries, but it is sensitive to noise. Larger values of $k$ smooth the decision boundary.
\end{remark}

\section{Evaluating on a Grid}

To visualize a classifier in two dimensions:
\begin{enumerate}
    \item Create a mesh of query points with resolution $h$.
    \item Apply the classifier to each grid point.
    \item Plot the predicted class regions.
    \item Overlay the training data.
\end{enumerate}

This produces a visualization of the decision boundary.

\begin{algorithm}[H]
\caption{$k$-Nearest Neighbors Classification}
\begin{algorithmic}[1]
\Require Training data $\{(x_i,C_i)\}_{i=1}^N$, query point $x$, number of neighbors $k$
\Ensure Predicted class label $\hat{C}(x)$
\State Compute distances $d_i \gets d(x,x_i)$ for $i=1,\dots,N$
\State Sort distances in increasing order
\State Let $\mathcal{N}_k(x)$ be the indices of the first $k$ points
\State Count class votes among $\{C_i : i \in \mathcal{N}_k(x)\}$
\State Return the class with the largest vote count
\end{algorithmic}
\end{algorithm}

\section{Meaningful Distance}

Nearest neighbor methods depend entirely on the distance function. If features have different units or scales, Euclidean distance can be dominated by variables with larger numerical ranges.

For example, suppose
\[
x=(\text{height in meters}, \text{income in dollars}).
\]
The income coordinate may dominate the distance unless features are scaled.

\begin{definition}[Standardization]
Given feature $j$, standardization transforms
\[
x_{ij}
\mapsto
\frac{x_{ij}-\bar{x}_j}{s_j},
\]
where $\bar{x}_j$ is the sample mean of feature $j$ and $s_j$ is its sample standard deviation.
\end{definition}

\begin{remark}
Centering and scaling are often essential preprocessing steps for nearest neighbor methods.
\end{remark}

\section{Curse of Dimensionality}

In high dimensions, nearest neighbor methods become less effective because distances become less informative.

\begin{remark}
Informally, in high dimensions, ``everybody is lonely.'' Points tend to be far apart, and the contrast between nearest and farthest neighbors decreases.
\end{remark}

\subsection{Volume Concentration}

Consider a unit ball in $\R^d$. The volume concentrates near the boundary as $d$ increases. For a point uniformly distributed in the unit ball, the probability that its radius is at most $r$ is
\[
\Prob(R\leq r)=r^d.
\]
Thus,
\[
\Prob(R>r)=1-r^d.
\]
For fixed $r<1$,
\[
r^d \to 0
\quad
\text{as}
\quad
d\to\infty.
\]
Therefore, most of the volume lies near the surface.

\begin{remark}
This explains why nearest neighbor methods may require very large sample sizes in high dimensions. The data become sparse relative to the feature space.
\end{remark}

\section{Implementation Notes}

\begin{lstlisting}[language=Python]
import numpy as np

def knn_predict(X_train, y_train, x_query, k=5):
    """
    Simple k-nearest neighbors classifier.
    X_train: array of shape (N, d)
    y_train: array of shape (N,)
    x_query: array of shape (d,)
    """
    distances = np.linalg.norm(X_train - x_query, axis=1)
    neighbor_idx = np.argsort(distances)[:k]
    neighbor_labels = y_train[neighbor_idx]

    labels, counts = np.unique(neighbor_labels, return_counts=True)
    return labels[np.argmax(counts)]
\end{lstlisting}

\begin{lstlisting}[language=Python]
def standardize_train_test(X_train, X_test):
    """
    Standardize using training-set statistics only.
    """
    mean = X_train.mean(axis=0)
    std = X_train.std(axis=0, ddof=1)

    X_train_scaled = (X_train - mean) / std
    X_test_scaled = (X_test - mean) / std

    return X_train_scaled, X_test_scaled
\end{lstlisting}

\chapter{Naive Bayes, LDA, and QDA}

\section{Naive Bayes Classifier}

\subsection{Bayesian Classification}

Let $\{C_k\}_{k=1}^K$ denote discrete classes and $x \in \R^d$ a feature vector.

By Bayes' rule,
\[
\Prob(C_k \mid x)
=
\frac{\pi_k \, \mathcal{L}_x(C_k)}{Z},
\]
where
\[
\pi_k = \Prob(C_k)
\quad \text{(prior)},
\quad
\mathcal{L}_x(C_k)=p(x \mid C_k)
\quad \text{(likelihood)},
\]
and $Z$ is a normalization constant independent of $k$.

\begin{proposition}[MAP Classification Rule]
The Bayes classifier assigns
\[
\hat{k}
=
\argmax_k \Prob(C_k \mid x)
=
\argmax_k \left[ \pi_k \, p(x \mid C_k) \right].
\]
\end{proposition}

\begin{remark}
The normalization constant cancels in the $\argmax$, so it need not be computed explicitly.
\end{remark}

\subsection{Naive Independence Assumption}

\begin{definition}[Naive Bayes Model]
Naive Bayes assumes conditional independence of features:
\[
p(x \mid C_k)
=
\prod_{\alpha=1}^d p(x_\alpha \mid C_k).
\]
\end{definition}

Thus, the classifier becomes
\[
\hat{k}
=
\argmax_k
\left[
\pi_k \prod_{\alpha=1}^d p(x_\alpha \mid C_k)
\right].
\]

\subsection{Gaussian Naive Bayes}

Assume each feature is conditionally Gaussian:
\[
p(x_\alpha \mid C_k)
=
\frac{1}{\sqrt{2\pi\sigma_{k,\alpha}^2}}
\exp\left(
-\frac{(x_\alpha-\mu_{k,\alpha})^2}{2\sigma_{k,\alpha}^2}
\right).
\]

\begin{definition}[Parameter Estimation]
For each class $k$ and feature $\alpha$, estimate
\[
\mu_{k,\alpha}
=
\frac{1}{N_k}\sum_{i:C_i=k} x_{i\alpha},
\quad
\sigma_{k,\alpha}^2
=
\frac{1}{N_k-1}\sum_{i:C_i=k}(x_{i\alpha}-\mu_{k,\alpha})^2.
\]
\end{definition}

Class priors are estimated as
\[
\pi_k
=
\frac{N_k}{N}.
\]

\subsection{Log-Domain Form}

To improve numerical stability, take logarithms:
\[
\hat{k}
=
\argmax_k
\left[
\log \pi_k
+
\sum_{\alpha=1}^d
\log p(x_\alpha \mid C_k)
\right].
\]

\begin{remark}
Working in the log-domain avoids numerical underflow from multiplying many small probabilities.
\end{remark}

\subsection{Advantages and Limitations}

\begin{proposition}[Scaling Invariance]
Naive Bayes is invariant to feature scaling (for Gaussian likelihoods), since each feature is normalized by its variance.
\end{proposition}

\begin{remark}
The independence assumption is often unrealistic. However, naive Bayes performs well in practice due to low variance and computational efficiency.
\end{remark}

\section{Discriminant Analysis}

Discriminant analysis assumes class-conditional Gaussian distributions:
\[
p(x \mid C_k)
=
\mathcal{N}(x;\mu_k,\Sigma_k).
\]

\subsection{Quadratic Discriminant Analysis (QDA)}

\begin{definition}[QDA Model]
Each class has its own covariance matrix:
\[
p(x \mid C_k)
=
\frac{1}{\sqrt{(2\pi)^d|\Sigma_k|}}
\exp\left(
-\frac{1}{2}(x-\mu_k)^T\Sigma_k^{-1}(x-\mu_k)
\right).
\]
\end{definition}

The discriminant function is obtained from the log-posterior:
\[
\delta_k(x)
=
\log \pi_k
-
\frac{1}{2}\log|\Sigma_k|
-
\frac{1}{2}(x-\mu_k)^T\Sigma_k^{-1}(x-\mu_k).
\]

\begin{proposition}[QDA Classification Rule]
\[
\hat{k}
=
\argmax_k \delta_k(x).
\]
\end{proposition}

\begin{remark}
The decision boundary between two classes is quadratic in $x$ due to the quadratic form $(x-\mu_k)^T\Sigma_k^{-1}(x-\mu_k)$.
\end{remark}

\subsection{Binary QDA Decision Boundary}

For two classes $C_1$ and $C_2$, compare
\[
(x-\mu_1)^T\Sigma_1^{-1}(x-\mu_1)+\log|\Sigma_1|
\]
and
\[
(x-\mu_2)^T\Sigma_2^{-1}(x-\mu_2)+\log|\Sigma_2|.
\]

The classification depends on which quantity is smaller (after including priors).

\section{Linear Discriminant Analysis (LDA)}

\begin{definition}[LDA Model]
LDA assumes a common covariance matrix across classes:
\[
\Sigma_k=\Sigma
\quad \forall k.
\]
\end{definition}

The discriminant function becomes
\[
\delta_k(x)
=
\log \pi_k
-
\frac{1}{2}(x-\mu_k)^T\Sigma^{-1}(x-\mu_k).
\]

\subsection{Linearization of the Decision Boundary}

Expand the quadratic term:
\[
(x-\mu_k)^T\Sigma^{-1}(x-\mu_k)
=
x^T\Sigma^{-1}x
-
2\mu_k^T\Sigma^{-1}x
+
\mu_k^T\Sigma^{-1}\mu_k.
\]

When comparing two classes, the $x^T\Sigma^{-1}x$ term cancels, yielding a linear function of $x$:
\[
\delta_k(x)
=
\mu_k^T\Sigma^{-1}x
-
\frac{1}{2}\mu_k^T\Sigma^{-1}\mu_k
+
\log \pi_k.
\]

\begin{theorem}[LDA Decision Boundary]
The decision boundary between classes is linear in $x$.
\end{theorem}

\begin{proof}
The discriminant difference between two classes is
\[
\delta_1(x)-\delta_2(x)
=
(\mu_1-\mu_2)^T\Sigma^{-1}x
-
\frac{1}{2}
\left(
\mu_1^T\Sigma^{-1}\mu_1
-
\mu_2^T\Sigma^{-1}\mu_2
\right)
+
\log \frac{\pi_1}{\pi_2}.
\]
This is an affine (linear + constant) function of $x$, so the boundary
\[
\delta_1(x)=\delta_2(x)
\]
is linear.
\end{proof}

\subsection{Parameter Estimation}

Means:
\[
\mu_k
=
\frac{1}{N_k}\sum_{i:C_i=k} x_i.
\]

Shared covariance:
\[
\Sigma
=
\frac{1}{N-K}
\sum_{k=1}^K
\sum_{i:C_i=k}
(x_i-\mu_k)(x_i-\mu_k)^T.
\]

\subsection{Comparison: LDA vs QDA}

\begin{itemize}
    \item QDA:
    \begin{itemize}
        \item Separate covariance matrices $\Sigma_k$
        \item Flexible (quadratic boundaries)
        \item High variance (many parameters)
    \end{itemize}
    \item LDA:
    \begin{itemize}
        \item Shared covariance $\Sigma$
        \item Linear boundaries
        \item Lower variance, better for small sample sizes
    \end{itemize}
\end{itemize}

\begin{remark}
The choice between LDA and QDA depends on whether class covariances are similar. If covariances differ significantly, QDA may perform better; otherwise LDA is more stable.
\end{remark}

\section{Algorithmic Summary}

\begin{algorithm}[H]
\caption{Gaussian Naive Bayes}
\begin{algorithmic}[1]
\Require Training data $\{(x_i,C_i)\}$, query point $x$
\Ensure Predicted class $\hat{k}$
\State Estimate $\mu_{k,\alpha}$ and $\sigma_{k,\alpha}^2$ for each class and feature
\State Estimate priors $\pi_k$
\For{each class $k$}
    \State Compute score:
    \[
    s_k = \log \pi_k + \sum_\alpha \log p(x_\alpha \mid C_k)
    \]
\EndFor
\State Return $\argmax_k s_k$
\end{algorithmic}
\end{algorithm}

\begin{algorithm}[H]
\caption{LDA Classification}
\begin{algorithmic}[1]
\Require Training data $\{(x_i,C_i)\}$, query point $x$
\Ensure Predicted class $\hat{k}$
\State Estimate $\mu_k$, shared $\Sigma$, and priors $\pi_k$
\For{each class $k$}
    \State Compute
    \[
    \delta_k(x)
    =
    \mu_k^T\Sigma^{-1}x
    -
    \frac{1}{2}\mu_k^T\Sigma^{-1}\mu_k
    +
    \log \pi_k
    \]
\EndFor
\State Return $\argmax_k \delta_k(x)$
\end{algorithmic}
\end{algorithm}

\begin{algorithm}[H]
\caption{QDA Classification}
\begin{algorithmic}[1]
\Require Training data $\{(x_i,C_i)\}$, query point $x$
\Ensure Predicted class $\hat{k}$
\State Estimate $\mu_k$, $\Sigma_k$, and priors $\pi_k$
\For{each class $k$}
    \State Compute
    \[
    \delta_k(x)
    =
    \log \pi_k
    -
    \frac{1}{2}\log|\Sigma_k|
    -
    \frac{1}{2}(x-\mu_k)^T\Sigma_k^{-1}(x-\mu_k)
    \]
\EndFor
\State Return $\argmax_k \delta_k(x)$
\end{algorithmic}
\end{algorithm}

\section{Implementation Notes}

\begin{lstlisting}[language=Python]
import numpy as np

def gaussian_naive_bayes_predict(X_train, y_train, x):
    classes = np.unique(y_train)
    scores = []

    for k in classes:
        X_k = X_train[y_train == k]
        mu = X_k.mean(axis=0)
        var = X_k.var(axis=0, ddof=1)

        log_prior = np.log(len(X_k) / len(X_train))
        log_likelihood = -0.5 * np.sum(
            np.log(2*np.pi*var) + ((x - mu)**2)/var
        )

        scores.append(log_prior + log_likelihood)

    return classes[np.argmax(scores)]
\end{lstlisting}

\begin{lstlisting}[language=Python]
def lda_predict(X_train, y_train, x):
    classes = np.unique(y_train)
    n, d = X_train.shape

    means = {}
    priors = {}

    for k in classes:
        X_k = X_train[y_train == k]
        means[k] = X_k.mean(axis=0)
        priors[k] = len(X_k) / n

    # shared covariance
    X_centered = []
    for k in classes:
        X_k = X_train[y_train == k]
        X_centered.append(X_k - means[k])
    X_centered = np.vstack(X_centered)

    Sigma = np.cov(X_centered, rowvar=False)
    Sigma_inv = np.linalg.inv(Sigma)

    scores = []
    for k in classes:
        mu = means[k]
        score = (
            mu @ Sigma_inv @ x
            - 0.5 * mu @ Sigma_inv @ mu
            + np.log(priors[k])
        )
        scores.append(score)

    return classes[np.argmax(scores)]
\end{lstlisting}

\chapter{Cross Validation}

\section{Model Evaluation}

Given a learning algorithm that produces an estimator $\hat{f}$ from training data, we seek to estimate its predictive performance on unseen data.

\begin{definition}[Generalization Error]
Let $(X,Y)$ be a new observation drawn from the same distribution as the training data. The generalization error is
\[
\mathcal{E}
=
\E\left[ L(Y,\hat{f}(X)) \right],
\]
where $L(\cdot,\cdot)$ is a loss function (e.g., squared error for regression or $0$-$1$ loss for classification).
\end{definition}

Since the true distribution is unknown, we approximate $\mathcal{E}$ using resampling methods.

\section{Train-Test Split and Complementary Subsets}

\begin{definition}[Hold-Out Method]
Split the data into two disjoint subsets:
\begin{itemize}
    \item Training set $\mathcal{D}_{\text{train}}$
    \item Test set $\mathcal{D}_{\text{test}}$
\end{itemize}
Train the model on $\mathcal{D}_{\text{train}}$ and evaluate on $\mathcal{D}_{\text{test}}$.
\end{definition}

\begin{remark}
The test error is
\[
\widehat{\mathcal{E}}
=
\frac{1}{|\mathcal{D}_{\text{test}}|}
\sum_{(x_i,y_i)\in \mathcal{D}_{\text{test}}}
L(y_i,\hat{f}(x_i)).
\]
\end{remark}

\begin{remark}
A single split can have high variance depending on how the data is partitioned. To reduce this variance, one may use repeated random splits (complementary subsets).
\end{remark}

\section{Repeated Random Subsampling}

\begin{definition}[Repeated Hold-Out]
Repeat the following procedure multiple times:
\begin{enumerate}
    \item Randomly split data into training and test sets.
    \item Train the model on the training set.
    \item Evaluate on the test set.
\end{enumerate}
Average the resulting test errors.
\end{definition}

\begin{remark}
This reduces variance relative to a single split but may reuse data unevenly across training and testing.
\end{remark}

\section{Leave-One-Out Cross Validation (LOOCV)}

\begin{definition}[LOOCV]
For a dataset of size $n$, perform $n$ training rounds. In each round:
\begin{itemize}
    \item Use $n-1$ observations for training
    \item Use the remaining observation for testing
\end{itemize}
\end{definition}

The LOOCV estimate of prediction error is
\[
\widehat{\mathcal{E}}_{\mathrm{LOOCV}}
=
\frac{1}{n}
\sum_{i=1}^n
L\big(y_i,\hat{f}^{(-i)}(x_i)\big),
\]
where $\hat{f}^{(-i)}$ is the model trained without the $i$-th observation.

\begin{remark}
LOOCV has low bias because nearly the full dataset is used for training, but it can be computationally expensive.
\end{remark}

\begin{remark}
For some models (e.g., linear regression), LOOCV can be computed efficiently without retraining using analytic formulas involving the hat matrix.
\end{remark}

\section{$k$-Fold Cross Validation}

\begin{definition}[$k$-Fold Cross Validation]
Partition the dataset into $k$ disjoint subsets (folds) of approximately equal size:
\[
\mathcal{D} = \mathcal{D}_1 \cup \cdots \cup \mathcal{D}_k.
\]
For each fold $j=1,\dots,k$:
\begin{itemize}
    \item Train on $\mathcal{D} \setminus \mathcal{D}_j$
    \item Test on $\mathcal{D}_j$
\end{itemize}
\end{definition}

The cross-validation estimate is
\[
\widehat{\mathcal{E}}_{\mathrm{CV}}
=
\frac{1}{k}
\sum_{j=1}^k
\frac{1}{|\mathcal{D}_j|}
\sum_{(x_i,y_i)\in \mathcal{D}_j}
L\big(y_i,\hat{f}^{(-j)}(x_i)\big).
\]

\begin{remark}
Common choices are $k=5$ or $k=10$. Smaller $k$ increases bias but reduces computational cost; larger $k$ reduces bias but increases variance and computation.
\end{remark}

\begin{proposition}
LOOCV is a special case of $k$-fold cross validation with $k=n$.
\end{proposition}

\section{Bias-Variance Tradeoff in Cross Validation}

\begin{itemize}
    \item LOOCV:
    \begin{itemize}
        \item Low bias
        \item High variance
        \item High computational cost
    \end{itemize}
    \item Small $k$ (e.g., $k=2$):
    \begin{itemize}
        \item Higher bias
        \item Lower variance
        \item Lower computational cost
    \end{itemize}
\end{itemize}

\section{Cross Validation Algorithm}

\begin{algorithm}[H]
\caption{$k$-Fold Cross Validation}
\begin{algorithmic}[1]
\Require Dataset $\mathcal{D}$, number of folds $k$
\Ensure Estimated prediction error $\widehat{\mathcal{E}}$
\State Partition $\mathcal{D}$ into $k$ folds $\{\mathcal{D}_j\}_{j=1}^k$
\For{$j=1$ to $k$}
    \State Train model $\hat{f}^{(-j)}$ on $\mathcal{D} \setminus \mathcal{D}_j$
    \State Compute error on $\mathcal{D}_j$
\EndFor
\State Average the errors across folds
\State \Return $\widehat{\mathcal{E}}$
\end{algorithmic}
\end{algorithm}

\section{Implementation Notes}

\begin{lstlisting}[language=Python]
from sklearn.model_selection import KFold
import numpy as np

def k_fold_cv(model, X, y, k=5):
    kf = KFold(n_splits=k, shuffle=True, random_state=0)
    errors = []

    for train_idx, test_idx in kf.split(X):
        X_train, X_test = X[train_idx], X[test_idx]
        y_train, y_test = y[train_idx], y[test_idx]

        model.fit(X_train, y_train)
        y_pred = model.predict(X_test)

        error = np.mean((y_test - y_pred)**2)
        errors.append(error)

    return np.mean(errors)
\end{lstlisting}

\begin{remark}
In practice, cross validation is used not only to estimate model performance but also to tune hyperparameters (e.g., $k$ in $k$-NN or $\lambda$ in regularization).
\end{remark}

\chapter{Decision Trees}

\section{Tree Structures}

\begin{definition}[Tree]
A tree is a connected acyclic graph with a hierarchical structure.
\end{definition}

\begin{definition}[Binary Tree]
A binary tree is a tree in which each node has at most two children:
\begin{itemize}
    \item A single \textbf{root} node
    \item Each node has at most a \textbf{left} and \textbf{right} child
    \item \textbf{Leaves} are terminal nodes with no children
\end{itemize}
\end{definition}

\begin{remark}
Trees are naturally defined recursively: each subtree is itself a tree.
\end{remark}

\subsection{$n$-ary Trees}

\begin{definition}[$n$-ary Tree]
An $n$-ary tree allows each node to have up to $n$ children.
\end{definition}

\section{Trees in Computation}

Trees appear in many applications:
\begin{itemize}
    \item Search structures: $k$d-trees, B-trees, R-trees, ball trees
    \item Decision making: classification and regression trees
\end{itemize}

\subsection{$k$d-Trees}

\begin{definition}[$k$d-Tree]
A $k$d-tree is a binary tree that recursively partitions $\R^d$ by splitting along coordinate axes.
\end{definition}

\begin{remark}
At each level, the split dimension cycles through coordinates, producing a space-partitioning structure.
\end{remark}

\section{Decision Trees}

\subsection{Recursive Partitioning}

A decision tree partitions the feature space recursively.

At a node with dataset $D$, we choose a split
\[
\theta=(j,t),
\]
where $j$ is a feature index and $t$ is a threshold.

The split defines two subsets:
\[
D_{\text{left}}(\theta)=\{x \in D : x_j \leq t\},
\quad
D_{\text{right}}(\theta)=\{x \in D : x_j > t\}.
\]

\subsection{Impurity Minimization}

We choose $\theta$ to minimize impurity:
\[
I(\theta)
=
\frac{n_{\text{left}}}{n}H(D_{\text{left}}(\theta))
+
\frac{n_{\text{right}}}{n}H(D_{\text{right}}(\theta)),
\]
where $H(\cdot)$ measures impurity and $n=|D|$.

\begin{definition}[Impurity Function]
An impurity function measures how mixed the class labels are in a dataset.
\end{definition}

\subsection{Gini Impurity}

\begin{definition}[Gini Impurity]
For a dataset $D$ with class proportions $p_1,\dots,p_K$,
\[
H(D)
=
\sum_{k=1}^K p_k(1-p_k).
\]
\end{definition}

\begin{remark}
Equivalently,
\[
H(D)=1-\sum_{k=1}^K p_k^2.
\]
\end{remark}

\subsection{Example: Root Impurity}

Suppose a dataset contains $8$ observations across $3$ classes with counts $(4,3,1)$. Then
\[
H
=
\frac{4}{8}\left(1-\frac{4}{8}\right)
+
\frac{3}{8}\left(1-\frac{3}{8}\right)
+
\frac{1}{8}\left(1-\frac{1}{8}\right).
\]
Compute each term:
\[
\frac{4}{8}\cdot\frac{4}{8}
=
\frac{16}{64},
\quad
\frac{3}{8}\cdot\frac{5}{8}
=
\frac{15}{64},
\quad
\frac{1}{8}\cdot\frac{7}{8}
=
\frac{7}{64}.
\]
Thus,
\[
H
=
\frac{16}{64}
+
\frac{15}{64}
+
\frac{7}{64}
=
\frac{38}{64}
=
\frac{19}{32}
\approx 0.59375.
\]

\subsection{Example: Split Impurity}

Suppose a split produces two partitions:

\paragraph{Left partition:} 2 samples, all from one class:
\[
H_{\text{left}}
=
\frac{2}{2}(1-\frac{2}{2})=0.
\]

\paragraph{Right partition:} 6 samples with counts $(3,2,1)$:
\[
H_{\text{right}}
=
\frac{3}{6}\left(1-\frac{3}{6}\right)
+
\frac{2}{6}\left(1-\frac{2}{6}\right)
+
\frac{1}{6}\left(1-\frac{1}{6}\right).
\]
Compute:
\[
=
\frac{3}{6}\cdot\frac{3}{6}
+
\frac{2}{6}\cdot\frac{4}{6}
+
\frac{1}{6}\cdot\frac{5}{6}
=
\frac{9}{36}
+
\frac{8}{36}
+
\frac{5}{36}
=
\frac{22}{36}
=
\frac{11}{18}
\approx 0.6111.
\]

\subsection{Weighted Impurity}

The total impurity after the split is
\[
I(\theta)
=
\frac{2}{8}\cdot 0
+
\frac{6}{8}\cdot \frac{11}{18}
=
\frac{6}{8}\cdot \frac{11}{18}
=
\frac{66}{144}
=
\frac{11}{24}.
\]

\begin{remark}
Impurity must be weighted by partition size. A perfectly pure small partition does not necessarily yield a good split overall.
\end{remark}

\section{Regression Trees}

For regression, impurity is measured by variance:
\[
H(D)
=
\frac{1}{|D|}
\sum_{x_i \in D}
(y_i-\bar{y})^2.
\]

\section{Greedy Tree Construction}

Decision trees are built greedily:
\begin{itemize}
    \item At each node, choose the best split minimizing impurity
    \item Recurse on resulting subsets
    \item Stop when a stopping criterion is met
\end{itemize}

\begin{algorithm}[H]
\caption{Decision Tree Training}
\begin{algorithmic}[1]
\Require Dataset $D$, stopping criterion
\Ensure Decision tree
\If{stopping criterion satisfied}
    \State Return leaf node with prediction
\EndIf
\For{each feature $j$ and threshold $t$}
    \State Compute split $\theta=(j,t)$
    \State Compute impurity $I(\theta)$
\EndFor
\State Choose $\theta^\star=\argmin_\theta I(\theta)$
\State Split data into $D_{\text{left}}, D_{\text{right}}$
\State Recursively build left and right subtrees
\State Return node with split $\theta^\star$
\end{algorithmic}
\end{algorithm}

\section{Properties of Decision Trees}

\begin{itemize}
    \item Nonparametric: no assumption on data distribution
    \item Capture nonlinear relationships
    \item Easy to interpret
    \item Can overfit without regularization
\end{itemize}

\section{Regularization and Practical Considerations}

Common techniques:
\begin{itemize}
    \item Maximum tree depth
    \item Minimum samples per leaf
    \item Pruning (post-training simplification)
\end{itemize}

\begin{remark}
Decision trees have low bias but high variance. Ensemble methods such as random forests reduce variance.
\end{remark}

\section{Implementation Notes}

\begin{lstlisting}[language=Python]
from sklearn.tree import DecisionTreeClassifier

model = DecisionTreeClassifier(max_depth=3)
model.fit(X_train, y_train)
y_pred = model.predict(X_test)
\end{lstlisting}

\begin{lstlisting}[language=Python]
from sklearn.tree import DecisionTreeRegressor

model = DecisionTreeRegressor(max_depth=3)
model.fit(X_train, y_train)
y_pred = model.predict(X_test)
\end{lstlisting}

\chapter{Scikit-Learn Pipeline and Random Forests}

\section{Scikit-Learn Pipeline and Grid Search}

\subsection{Pipelines}

\begin{definition}[Pipeline]
A pipeline chains multiple transformations and a final estimator into a single object:
\[
X \rightarrow \text{Transform}_1 \rightarrow \cdots \rightarrow \text{Transform}_m \rightarrow \hat{f}.
\]
\end{definition}

\begin{remark}
Pipelines ensure that preprocessing steps (e.g., scaling, PCA) are applied consistently during training and testing.
\end{remark}

\begin{example}[Pipeline Structure]
\begin{lstlisting}[language=Python]
from sklearn.pipeline import Pipeline
from sklearn.preprocessing import StandardScaler
from sklearn.decomposition import PCA
from sklearn.linear_model import LogisticRegression

pipe = Pipeline([
    ("scaler", StandardScaler()),
    ("pca", PCA(n_components=5)),
    ("model", LogisticRegression())
])
\end{lstlisting}
\end{example}

\subsection{Grid Search with Cross Validation}

\begin{definition}[Grid Search]
Grid search selects hyperparameters by evaluating a model over a predefined parameter grid using cross validation.
\end{definition}

\begin{example}[Grid Search with Pipeline]
\begin{lstlisting}[language=Python]
from sklearn.model_selection import GridSearchCV

param_grid = {
    "pca__n_components": [2, 5, 10],
    "model__C": [0.1, 1, 10]
}

grid = GridSearchCV(pipe, param_grid, cv=5)
grid.fit(X_train, y_train)

best_model = grid.best_estimator_
\end{lstlisting}
\end{example}

\begin{remark}
The notation \texttt{step\_\_param} specifies a parameter for a given pipeline step.
\end{remark}

\begin{remark}
Grid search combined with cross validation provides a principled way to tune hyperparameters while avoiding overfitting to a single train-test split.
\end{remark}

\section{Random Forests}

\subsection{Random Trees}

\begin{definition}[Randomized Tree]
A randomized decision tree selects a random subset of features when searching for the best split at each node.
\end{definition}

\begin{remark}
In high dimensions, searching all features is computationally expensive. Random feature selection reduces cost and introduces diversity across trees.
\end{remark}

\subsection{Forest of Random Trees}

\begin{definition}[Random Forest]
A random forest is an ensemble of randomized decision trees:
\[
\{\hat{f}_1,\dots,\hat{f}_B\}.
\]
\end{definition}

Predictions are combined:
\begin{itemize}
    \item Classification: majority vote
    \item Regression: average prediction
\end{itemize}

\begin{definition}[Bootstrap Sampling]
Each tree is trained on a bootstrap sample of the data, obtained by sampling with replacement.
\end{definition}

\subsection{Bagging and Variance Reduction}

\begin{proposition}[Variance Reduction via Averaging]
Let $\hat{f}_1,\dots,\hat{f}_B$ be independent estimators with variance $\sigma^2$. Then the average
\[
\bar{f}=\frac{1}{B}\sum_{b=1}^B \hat{f}_b
\]
has variance
\[
\Var(\bar{f})=\frac{\sigma^2}{B}.
\]
\end{proposition}

\begin{proof}
Using independence,
\[
\Var\left(\frac{1}{B}\sum_{b=1}^B \hat{f}_b\right)
=
\frac{1}{B^2}\sum_{b=1}^B \Var(\hat{f}_b)
=
\frac{B\sigma^2}{B^2}
=
\frac{\sigma^2}{B}.
\]
\end{proof}

\begin{remark}
In practice, trees are not independent, but randomization reduces correlation, improving variance reduction.
\end{remark}

\subsection{Connection to the Central Limit Theorem}

\begin{remark}
As the number of trees $B$ increases, the average prediction stabilizes. Under mild conditions, the distribution of the ensemble prediction approaches a normal distribution due to averaging effects.
\end{remark}

\section{Feature Selection in Random Forests}

\begin{definition}[Feature Importance]
Feature importance measures how often and how effectively a feature is used to split the data.
\end{definition}

A common measure is the total impurity reduction contributed by a feature across all trees.

\begin{remark}
Features that frequently produce large impurity reductions are considered more important.
\end{remark}

\section{Assumptions and Limitations}

\begin{itemize}
    \item Decision trees produce axis-aligned splits, which may be suboptimal for rotated or highly correlated data
    \item No explicit distance function is required
    \item Random forests reduce overfitting relative to single trees
\end{itemize}

\section{Divide and Conquer Perspective}

Decision trees partition the feature space recursively:
\[
\R^d = \bigcup_{m=1}^M R_m,
\]
where each region $R_m$ corresponds to a leaf.

For regression, predictions are piecewise constant:
\[
\hat{f}(x)
=
\sum_{m=1}^M c_m \ind(x \in R_m),
\]
where
\[
c_m
=
\frac{1}{|R_m|}
\sum_{x_i \in R_m} y_i.
\]

\begin{remark}
This can be interpreted as minimizing within-region variance:
\[
\sum_{m=1}^M \sum_{x_i \in R_m} (y_i - c_m)^2.
\]
\end{remark}

\section{Random Forest Algorithm}

\begin{algorithm}[H]
\caption{Random Forest Training}
\begin{algorithmic}[1]
\Require Training data $\mathcal{D}$, number of trees $B$
\Ensure Ensemble of trees $\{\hat{f}_b\}_{b=1}^B$
\For{$b=1$ to $B$}
    \State Draw bootstrap sample $\mathcal{D}_b$ from $\mathcal{D}$
    \State Train decision tree $\hat{f}_b$:
    \begin{itemize}
        \item At each split, select a random subset of features
        \item Choose the best split among those features
    \end{itemize}
\EndFor
\State \Return $\{\hat{f}_b\}_{b=1}^B$
\end{algorithmic}
\end{algorithm}

\begin{algorithm}[H]
\caption{Random Forest Prediction}
\begin{algorithmic}[1]
\Require Trees $\{\hat{f}_b\}_{b=1}^B$, query point $x$
\Ensure Prediction $\hat{y}$
\If{classification}
    \State $\hat{y} \gets$ majority vote of $\{\hat{f}_b(x)\}$
\Else
    \State $\hat{y} \gets \frac{1}{B}\sum_{b=1}^B \hat{f}_b(x)$
\EndIf
\State \Return $\hat{y}$
\end{algorithmic}
\end{algorithm}

\section{Implementation Notes}

\begin{lstlisting}[language=Python]
from sklearn.pipeline import Pipeline
from sklearn.model_selection import GridSearchCV
from sklearn.ensemble import RandomForestClassifier
from sklearn.preprocessing import StandardScaler

pipe = Pipeline([
    ("scaler", StandardScaler()),
    ("rf", RandomForestClassifier())
])

param_grid = {
    "rf__n_estimators": [100, 200],
    "rf__max_depth": [None, 5, 10],
    "rf__max_features": ["sqrt", "log2"]
}

grid = GridSearchCV(pipe, param_grid, cv=5)
grid.fit(X_train, y_train)

best_model = grid.best_estimator_
\end{lstlisting}

\begin{remark}
Random forests typically require minimal preprocessing and perform well out-of-the-box, making them a strong baseline model.
\end{remark}

\chapter{Logistic Regression}

\section{Binary Classification and the Logistic Function}

In binary classification, we model the posterior probability
\[
\pi(x) = \Prob(C_1 \mid x).
\]

\begin{definition}[Logistic (Sigmoid) Function]
The logistic function is
\[
\sigma(t)
=
\frac{1}{1+e^{-t}}.
\]
\end{definition}

\begin{remark}
The logistic function maps $\R \to (0,1)$, making it suitable for modeling probabilities.
\end{remark}

\section{Log-Odds and Linear Model}

Define the log-odds (logit) transformation:
\[
\log \frac{\Prob(C_1 \mid x)}{\Prob(C_0 \mid x)}
=
\log \frac{\pi(x)}{1-\pi(x)}.
\]

\begin{definition}[Logistic Regression Model]
Assume the log-odds are linear in the features:
\[
\log \frac{\pi(x)}{1-\pi(x)}
=
\beta_0 + \beta^T x.
\]
\end{definition}

Solving for $\pi(x)$,
\[
\pi(x)
=
\frac{1}{1+e^{-(\beta_0+\beta^T x)}}.
\]

\begin{remark}
This is equivalent to applying the sigmoid function to a linear predictor:
\[
\pi(x) = \sigma(\beta_0+\beta^T x).
\]
\end{remark}

\section{Likelihood and Estimation}

Given training data
\[
\{(x_i,c_i)\}_{i=1}^n,
\quad
c_i \in \{0,1\},
\]
the likelihood is
\[
\mathcal{L}(\beta)
=
\prod_{i=1}^n
\pi(x_i;\beta)^{c_i}
\left[1-\pi(x_i;\beta)\right]^{1-c_i}.
\]

\begin{definition}[Log-Likelihood]
The log-likelihood is
\[
\ell(\beta)
=
\sum_{i=1}^n
\left[
c_i \log \pi(x_i;\beta)
+
(1-c_i)\log(1-\pi(x_i;\beta))
\right].
\]
\end{definition}

\begin{proposition}[Maximum Likelihood Estimation]
The logistic regression estimator is
\[
\hat{\beta}
=
\argmax_{\beta} \ell(\beta).
\]
\end{proposition}

\subsection{Gradient and Optimization}

Let $\pi_i = \pi(x_i;\beta)$. Then
\[
\frac{\partial \ell}{\partial \beta}
=
\sum_{i=1}^n (c_i - \pi_i)x_i.
\]

\begin{proof}
We use
\[
\frac{d}{dt}\sigma(t) = \sigma(t)(1-\sigma(t)).
\]
Differentiating the log-likelihood,
\[
\frac{\partial \ell}{\partial \beta}
=
\sum_{i=1}^n
\left[
\frac{c_i}{\pi_i}\frac{\partial \pi_i}{\partial \beta}
-
\frac{1-c_i}{1-\pi_i}\frac{\partial \pi_i}{\partial \beta}
\right].
\]
Using
\[
\frac{\partial \pi_i}{\partial \beta}
=
\pi_i(1-\pi_i)x_i,
\]
we obtain
\[
\frac{\partial \ell}{\partial \beta}
=
\sum_{i=1}^n
(c_i-\pi_i)x_i.
\]
\end{proof}

\begin{remark}
The log-likelihood is concave, so standard optimization methods (gradient ascent, Newton's method) converge to the global optimum.
\end{remark}

\section{Prediction}

\begin{definition}[Prediction Rule]
Given $\hat{\beta}$, predict
\[
\hat{C}(x)
=
\begin{cases}
1, & \pi(x;\hat{\beta}) \geq 0.5,\\
0, & \text{otherwise}.
\end{cases}
\]
\end{definition}

\section{Decision Boundary}

The decision boundary is given by
\[
\beta_0 + \beta^T x = 0,
\]
which is linear in $x$.

\begin{remark}
Thus, logistic regression produces linear decision boundaries, similar to LDA.
\end{remark}

\section{Multiclass Logistic Regression}

For $K$ classes, define
\[
\pi_k(x)
=
\Prob(C_k \mid x).
\]

\begin{definition}[Softmax Model]
\[
\pi_k(x)
=
\frac{e^{\beta_k^T x}}{\sum_{j=1}^K e^{\beta_j^T x}}.
\]
\end{definition}

The log-likelihood is
\[
\ell(\beta)
=
\sum_{i=1}^n
\log \pi_{c_i}(x_i;\beta).
\]

\begin{remark}
This is known as multinomial logistic regression or softmax regression.
\end{remark}

\section{Discriminative vs Generative Models}

\begin{definition}[Discriminative Model]
A discriminative model directly estimates
\[
\Prob(C_k \mid x).
\]
\end{definition}

\begin{definition}[Generative Model]
A generative model estimates
\[
p(x \mid C_k)
\quad \text{and} \quad \Prob(C_k),
\]
and uses Bayes' rule to compute posteriors.
\end{definition}

\begin{remark}
Logistic regression is discriminative, whereas naive Bayes and LDA/QDA are generative.
\end{remark}

\section{Regularization}

To prevent overfitting, add a penalty:

\subsection{L2 Regularization (Ridge)}

\[
\hat{\beta}
=
\argmax_{\beta}
\left[
\ell(\beta) - \lambda \norm{\beta}_2^2
\right].
\]

\subsection{L1 Regularization (Lasso)}

\[
\hat{\beta}
=
\argmax_{\beta}
\left[
\ell(\beta) - \lambda \norm{\beta}_1
\right].
\]

\section{Algorithm}

\begin{algorithm}[H]
\caption{Logistic Regression via Gradient Ascent}
\begin{algorithmic}[1]
\Require Training data $\{(x_i,c_i)\}$, learning rate $\eta$
\Ensure Parameter estimate $\hat{\beta}$
\State Initialize $\beta$
\For{$t=1$ to max iterations}
    \State Compute $\pi_i = \sigma(\beta^T x_i)$ for all $i$
    \State Compute gradient:
    \[
    g = \sum_i (c_i - \pi_i)x_i
    \]
    \State Update:
    \[
    \beta \gets \beta + \eta g
    \]
\EndFor
\State \Return $\beta$
\end{algorithmic}
\end{algorithm}

\section{Implementation Notes}

\begin{lstlisting}[language=Python]
import numpy as np

def sigmoid(z):
    return 1 / (1 + np.exp(-z))

def logistic_regression_fit(X, y, lr=0.01, n_iter=1000):
    n, d = X.shape
    beta = np.zeros(d)

    for _ in range(n_iter):
        z = X @ beta
        p = sigmoid(z)
        grad = X.T @ (y - p)
        beta += lr * grad

    return beta

def logistic_regression_predict(X, beta):
    probs = sigmoid(X @ beta)
    return (probs >= 0.5).astype(int)
\end{lstlisting}

\chapter{Support Vector Machines}

\section{Maximum Margin Classification}

Support Vector Machines (SVMs) construct a separating hyperplane between classes with maximum margin.

\begin{definition}[Hyperplane]
A hyperplane in $\R^d$ is defined as
\[
w^T x - b = 0,
\]
where $w \in \R^d$ is the normal vector and $b \in \R$ is the bias.
\end{definition}

\begin{remark}
The sign of $w^T x - b$ determines the predicted class.
\end{remark}

\section{Hard Margin SVM}

Assume the data are linearly separable with labels $y_i \in \{-1,1\}$.

We define two parallel hyperplanes:
\[
w^T x - b = 1,
\quad
w^T x - b = -1.
\]

\subsection{Margin Width}

\begin{proposition}[Margin Width]
The distance between the two margin hyperplanes is
\[
\frac{2}{\norm{w}_2}.
\]
\end{proposition}

\begin{proof}
Let $x_0$ satisfy
\[
w^T x_0 - b = -1.
\]
Let $x_1 = x_0 + d\frac{w}{\norm{w}}$ lie on the other hyperplane:
\[
w^T x_1 - b = 1.
\]
Substitute:
\[
w^T\left(x_0 + d\frac{w}{\norm{w}}\right) - b
=
w^T x_0 - b + d \frac{w^T w}{\norm{w}}.
\]
Since $w^T x_0 - b = -1$ and $w^T w = \norm{w}^2$,
\[
-1 + d \frac{\norm{w}^2}{\norm{w}} = 1.
\]
Thus,
\[
-1 + d \norm{w} = 1
\quad \Rightarrow \quad
d = \frac{2}{\norm{w}}.
\]
\end{proof}

\subsection{Optimization Problem}

Maximizing the margin is equivalent to minimizing $\norm{w}_2^2$:

\begin{definition}[Hard Margin SVM]
\[
\min_{w,b} \norm{w}_2^2
\quad \text{s.t.} \quad
y_i(w^T x_i - b) \ge 1
\quad \forall i.
\]
\end{definition}

\begin{remark}
This is a convex quadratic optimization problem with linear constraints.
\end{remark}

\section{Support Vectors}

\begin{definition}[Support Vectors]
The support vectors are the data points that lie exactly on the margin:
\[
y_i(w^T x_i - b) = 1.
\]
\end{definition}

\begin{remark}
Only support vectors determine the optimal hyperplane. Points farther from the boundary do not affect the solution.
\end{remark}

\section{Soft Margin SVM}

When classes are not linearly separable, introduce slack variables $\xi_i \ge 0$.

\begin{definition}[Soft Margin SVM]
\[
\min_{w,b,\xi}
\quad
\norm{w}_2^2 + C \sum_{i=1}^n \xi_i
\]
subject to
\[
y_i(w^T x_i - b) \ge 1 - \xi_i,
\quad
\xi_i \ge 0.
\]
\end{definition}

\begin{remark}
The parameter $C > 0$ controls the tradeoff between margin size and classification errors.
\end{remark}

\subsection{Hinge Loss Formulation}

The soft-margin objective can be written as
\[
\min_{w,b}
\quad
\norm{w}_2^2
+
C \sum_{i=1}^n
\max\{0,\,1 - y_i(w^T x_i - b)\}.
\]

\begin{remark}
The hinge loss penalizes points that are inside the margin or misclassified.
\end{remark}

\section{Kernel Methods}

\subsection{Feature Transformation}

Instead of working in the original space, map data to a higher-dimensional feature space:
\[
x \mapsto \phi(x).
\]

The SVM optimization then uses
\[
w^T \phi(x).
\]

\subsection{Kernel Trick}

\begin{definition}[Kernel Function]
A kernel function is defined as
\[
k(x,x') = \phi(x)^T \phi(x').
\]
\end{definition}

\begin{remark}
The kernel trick allows computation in high-dimensional feature spaces without explicitly computing $\phi(x)$.
\end{remark}

\subsection{Common Kernels}

\begin{itemize}
    \item Polynomial kernel:
    \[
    k(x,x') = (x^T x' + c)^d
    \]
    \item Radial basis function (RBF):
    \[
    k(x,x') = \exp\left(-\frac{\norm{x-x'}^2}{2\sigma^2}\right)
    \]
\end{itemize}

\begin{remark}
RBF kernels produce nonlinear decision boundaries by mapping data into infinite-dimensional feature spaces.
\end{remark}

\section{Dual Formulation (Key Insight)}

The SVM solution can be expressed in terms of inner products:
\[
w = \sum_{i=1}^n \alpha_i y_i x_i.
\]

Thus predictions become
\[
f(x) = \sum_{i=1}^n \alpha_i y_i k(x_i,x) - b.
\]

\begin{remark}
Only support vectors have nonzero $\alpha_i$, making the model sparse.
\end{remark}

\section{Algorithmic Summary}

\begin{algorithm}[H]
\caption{Support Vector Machine (Conceptual)}
\begin{algorithmic}[1]
\Require Training data $\{(x_i,y_i)\}$
\Ensure Classifier $(w,b)$
\State Solve quadratic optimization problem
\State Identify support vectors
\State Construct decision function:
\[
f(x) = \sum_i \alpha_i y_i k(x_i,x) - b
\]
\State \Return classifier
\end{algorithmic}
\end{algorithm}

\section{Properties and Comparison}

\begin{itemize}
    \item Maximizes geometric margin
    \item Robust to high-dimensional data
    \item Depends only on support vectors
    \item Kernel trick enables nonlinear classification
\end{itemize}

\begin{remark}
Unlike nearest neighbors, SVM does not rely on a distance metric in the original space, but rather on inner products (or kernels).
\end{remark}

\section{Implementation Notes}

\begin{lstlisting}[language=Python]
from sklearn.svm import SVC

model = SVC(kernel="rbf", C=1.0, gamma="scale")
model.fit(X_train, y_train)

y_pred = model.predict(X_test)
\end{lstlisting}

\begin{lstlisting}[language=Python]
from sklearn.pipeline import Pipeline
from sklearn.preprocessing import StandardScaler
from sklearn.svm import SVC

pipe = Pipeline([
    ("scaler", StandardScaler()),
    ("svm", SVC())
])

pipe.fit(X_train, y_train)
\end{lstlisting}

\begin{remark}
SVMs are sensitive to feature scaling, so standardization is typically required.
\end{remark}

\chapter{$k$-Means Clustering}

\section{Clustering}

Clustering is the task of grouping data points into subsets (clusters) such that points in the same cluster are more similar to each other than to points in other clusters.

\begin{itemize}
    \item Discover species based on features
    \item Segment images by pixel similarity
    \item Group documents by topic
\end{itemize}

\section{Types of Clustering Algorithms}

\subsection{Flat (Partitioning) Methods}

\begin{itemize}
    \item Start with an initial partition of the data
    \item Iteratively refine cluster assignments
\end{itemize}

\subsection{Hierarchical Methods}

\begin{itemize}
    \item \textbf{Agglomerative (bottom-up)}: merge closest clusters
    \item \textbf{Divisive (top-down)}: split clusters recursively
\end{itemize}

\section{$k$-Means Clustering}

\subsection{Optimization Problem}

Given data $\{x_l\}_{l=1}^n$ with $x_l \in \R^d$, partition the data into $k$ clusters $\{C_i\}_{i=1}^k$.

\begin{definition}[$k$-Means Objective]
\[
\hat{C}
=
\argmin_{C}
\sum_{i=1}^k
\sum_{x \in C_i}
\norm{x-\mu_i}_2^2,
\]
where
\[
\mu_i
=
\frac{1}{|C_i|}
\sum_{x \in C_i} x.
\]
\end{definition}

\begin{remark}
This objective is also called the within-cluster sum of squares (WCSS).
\end{remark}

\subsection{Optimality of the Mean}

\begin{proposition}
Given a fixed cluster $C$, the point minimizing
\[
\sum_{x \in C} \norm{x - \mu}_2^2
\]
is the mean
\[
\mu = \frac{1}{|C|}\sum_{x \in C} x.
\]
\end{proposition}

\begin{proof}
Let
\[
f(\mu)
=
\sum_{x \in C} \norm{x-\mu}_2^2.
\]
Expanding,
\[
f(\mu)
=
\sum_{x \in C} (x-\mu)^T(x-\mu).
\]
Taking gradient,
\[
\grad_\mu f(\mu)
=
-2\sum_{x \in C} (x-\mu).
\]
Setting equal to zero,
\[
\sum_{x \in C} (x-\mu)=0
\quad \Rightarrow \quad
\mu=\frac{1}{|C|}\sum_{x \in C} x.
\]
\end{proof}

\section{Algorithm}

$k$-means is an iterative algorithm alternating between assignment and update steps.

\begin{algorithm}[H]
\caption{$k$-Means Clustering}
\begin{algorithmic}[1]
\Require Data $\{x_i\}_{i=1}^n$, number of clusters $k$
\Ensure Clusters $\{C_i\}$ and centroids $\{\mu_i\}$
\State Initialize centroids $\mu_1,\dots,\mu_k$
\Repeat
    \State \textbf{Assignment step:}
    \For{each data point $x_i$}
        \State Assign $x_i$ to cluster
        \[
        C_j = \argmin_{j} \norm{x_i-\mu_j}_2
        \]
    \EndFor
    \State \textbf{Update step:}
    \For{each cluster $C_j$}
        \State Update centroid
        \[
        \mu_j = \frac{1}{|C_j|}\sum_{x \in C_j} x
        \]
    \EndFor
\Until{convergence}
\end{algorithmic}
\end{algorithm}

\begin{remark}
Each iteration decreases the objective function, so the algorithm converges to a local minimum.
\end{remark}

\section{Inertia}

\begin{definition}[Inertia]
Inertia is the total within-cluster sum of squares:
\[
\text{Inertia}
=
\sum_{i=1}^k
\sum_{x \in C_i}
\norm{x-\mu_i}_2^2.
\]
\end{definition}

\begin{remark}
\begin{itemize}
    \item Lower inertia indicates tighter clusters
    \item Inertia always decreases as $k$ increases
\end{itemize}
\end{remark}

\section{Choosing the Number of Clusters}

\subsection{Elbow Method}

Plot inertia as a function of $k$ and look for a point where the decrease slows significantly.

\begin{remark}
The ``elbow'' corresponds to a good tradeoff between model complexity and fit.
\end{remark}

\section{Limitations}

\begin{itemize}
    \item Sensitive to initialization
    \item Converges to local minima
    \item Assumes spherical clusters (Euclidean distance)
    \item Sensitive to feature scaling
\end{itemize}

\subsection{Initialization Strategies}

\begin{itemize}
    \item Random initialization
    \item Multiple restarts (parameter \texttt{n\_init})
    \item $k$-means++ initialization
\end{itemize}

\section{Variants}

\subsection{$k$-Medians}

\begin{definition}[$k$-Medians Objective]
\[
\min_{C}
\sum_{i=1}^k
\sum_{x \in C_i}
\norm{x-\mu_i}_1,
\]
where $\mu_i$ is the coordinate-wise median.
\end{definition}

\begin{remark}
Using $L_1$ distance makes the method more robust to outliers.
\end{remark}

\section{Geometric Interpretation}

$k$-means partitions the space into Voronoi regions:
\[
R_j = \{x : \norm{x-\mu_j}_2 \leq \norm{x-\mu_l}_2 \ \forall l\}.
\]

\begin{remark}
The decision boundaries between clusters are linear (perpendicular bisectors).
\end{remark}

\section{Implementation Notes}

\begin{lstlisting}[language=Python]
from sklearn.cluster import KMeans

model = KMeans(n_clusters=3, n_init=10)
model.fit(X)

labels = model.labels_
centroids = model.cluster_centers_
\end{lstlisting}

\begin{lstlisting}[language=Python]
import matplotlib.pyplot as plt

inertia = []
k_values = range(1, 10)

for k in k_values:
    model = KMeans(n_clusters=k, n_init=10)
    model.fit(X)
    inertia.append(model.inertia_)

plt.plot(k_values, inertia)
plt.xlabel("k")
plt.ylabel("Inertia")
plt.title("Elbow Method")
plt.show()
\end{lstlisting}

\chapter{Gaussian Mixture Models}

\section{Probabilistic Clustering}

Gaussian Mixture Models (GMMs) provide a probabilistic approach to clustering by modeling the data as a mixture of Gaussian distributions.

\begin{definition}[Gaussian Mixture Model]
A GMM with $K$ components assumes
\[
p(x)
=
\sum_{k=1}^K \pi_k \, \mathcal{N}(x;\mu_k,\Sigma_k),
\]
where
\begin{itemize}
    \item $\pi_k \ge 0$, $\sum_{k=1}^K \pi_k = 1$ (mixture weights)
    \item $\mu_k \in \R^d$ (means)
    \item $\Sigma_k \in \R^{d\times d}$ (covariances)
\end{itemize}
\end{definition}

\begin{remark}
Each component represents a cluster, but membership is probabilistic rather than deterministic.
\end{remark}

\section{Latent Variables}

Introduce latent variables $z_{ki}$:
\[
z_{ki} =
\begin{cases}
1, & \text{if } x_i \text{ belongs to component } k,\\
0, & \text{otherwise}.
\end{cases}
\]

Since cluster memberships are unknown, we treat $z_{ki}$ as hidden variables.

\section{Likelihood Function}

Given data $\{x_i\}_{i=1}^n$, the likelihood is
\[
L(\theta; x)
=
\prod_{i=1}^n
\left[
\sum_{k=1}^K \pi_k \, \mathcal{N}(x_i;\mu_k,\Sigma_k)
\right].
\]

\begin{remark}
The log-likelihood involves a log of sums, making direct maximization difficult.
\end{remark}

\section{Expectation-Maximization (EM) Algorithm}

The EM algorithm iteratively maximizes the likelihood in the presence of latent variables.

\begin{definition}[EM Framework]
Given current parameters $\theta$, compute updated parameters $\theta'$ such that
\[
p(D \mid \theta') \ge p(D \mid \theta).
\]
\end{definition}

\subsection{E-Step}

Compute posterior probabilities (responsibilities):
\[
\gamma_{ik}
=
\Prob(z_{ki}=1 \mid x_i)
=
\frac{\pi_k \, \mathcal{N}(x_i;\mu_k,\Sigma_k)}
{\sum_{j=1}^K \pi_j \, \mathcal{N}(x_i;\mu_j,\Sigma_j)}.
\]

\begin{remark}
$\gamma_{ik}$ represents the degree to which point $x_i$ belongs to cluster $k$.
\end{remark}

\subsection{M-Step}

Update parameters using responsibilities:

\[
N_k = \sum_{i=1}^n \gamma_{ik}.
\]

\[
\pi_k = \frac{N_k}{n}.
\]

\[
\mu_k = \frac{1}{N_k} \sum_{i=1}^n \gamma_{ik} x_i.
\]

\[
\Sigma_k = \frac{1}{N_k}
\sum_{i=1}^n
\gamma_{ik}
(x_i-\mu_k)(x_i-\mu_k)^T.
\]

\begin{remark}
The M-step computes weighted means and covariances, where weights are the responsibilities.
\end{remark}

\subsection{Algorithm}

\begin{algorithm}[H]
\caption{Gaussian Mixture Model via EM}
\begin{algorithmic}[1]
\Require Data $\{x_i\}$, number of components $K$
\Ensure Parameters $\{\pi_k,\mu_k,\Sigma_k\}$
\State Initialize parameters
\Repeat
    \State \textbf{E-step:} Compute $\gamma_{ik}$ for all $i,k$
    \State \textbf{M-step:} Update $\pi_k,\mu_k,\Sigma_k$
\Until{convergence}
\end{algorithmic}
\end{algorithm}

\section{Connection to $k$-Means}

\begin{remark}
$k$-means can be viewed as a limiting case of GMMs:
\begin{itemize}
    \item Equal covariances $\Sigma_k = \sigma^2 I$
    \item $\sigma^2 \to 0$
\end{itemize}
This leads to hard assignments instead of probabilistic memberships.
\end{remark}

\section{Soft Clustering}

\begin{definition}[Soft Clustering]
In GMMs, each data point belongs to each cluster with probability $\gamma_{ik}$.
\end{definition}

\begin{remark}
This contrasts with $k$-means, where each point is assigned to exactly one cluster.
\end{remark}

\section{Properties and Limitations}

\begin{itemize}
    \item Captures elliptical clusters via covariance matrices
    \item Provides probabilistic cluster assignments
    \item More flexible than $k$-means
    \item Sensitive to initialization
    \item May converge to local optima
\end{itemize}

\section{Model Selection}

Choosing $K$:
\begin{itemize}
    \item Likelihood-based methods
    \item Information criteria (AIC, BIC)
\end{itemize}

\begin{remark}
Unlike inertia in $k$-means, likelihood does not necessarily monotonically improve with $K$ when penalization is used.
\end{remark}

\section{Comparison with $k$-Means}

\begin{itemize}
    \item GMM:
    \begin{itemize}
        \item Probabilistic model
        \item Elliptical clusters
        \item Soft assignments
    \end{itemize}
    \item $k$-means:
    \begin{itemize}
        \item Deterministic assignments
        \item Spherical clusters
        \item Faster and simpler
    \end{itemize}
\end{itemize}

\section{Implementation Notes}

\begin{lstlisting}[language=Python]
from sklearn.mixture import GaussianMixture

model = GaussianMixture(n_components=3)
model.fit(X)

labels = model.predict(X)
probs = model.predict_proba(X)
\end{lstlisting}

\begin{remark}
\texttt{predict\_proba} returns the soft cluster assignments $\gamma_{ik}$.
\end{remark}

\section{Further Clustering Methods}

\begin{itemize}
    \item Agglomerative clustering (hierarchical)
    \item DBSCAN (density-based clustering)
\end{itemize}

\begin{remark}
These methods do not require specifying the number of clusters in advance and can capture more complex cluster shapes.
\end{remark}

\chapter{Spectral Embedding and Spectral Clustering}

\section{Spectral Methods}

Spectral methods construct low-dimensional representations of data using eigenvectors of matrices derived from pairwise similarities.

\begin{definition}[Spectral Embedding]
Spectral embedding maps data points into a lower-dimensional space using eigenvectors of a graph-based matrix constructed from similarities.
\end{definition}

\begin{definition}[Spectral Clustering]
Spectral clustering applies a standard clustering method (e.g., $k$-means) to the embedded coordinates obtained from spectral embedding.
\end{definition}

\section{Graphs}

\begin{definition}[Graph]
A graph is defined as
\[
G=(V,E),
\]
where $V$ is a set of vertices and $E$ is a set of edges.
\end{definition}

\begin{definition}[Vertex and Edge]
\begin{itemize}
    \item A \textbf{vertex} is a node in the graph
    \item An \textbf{edge} connects two vertices
\end{itemize}
\end{definition}

\section{Similarity Graphs}

\begin{definition}[Similarity Graph]
A similarity graph connects data points based on a notion of similarity.
\end{definition}

Common constructions:
\begin{itemize}
    \item $\varepsilon$-neighborhood graph
    \item $k$-nearest neighbor graph
    \item Fully connected graph with weighted edges
\end{itemize}

\begin{remark}
Edges may be weighted by similarity values such as Gaussian kernels:
\[
w_{ij} = \exp\left(-\frac{\norm{x_i - x_j}^2}{2\sigma^2}\right).
\]
\end{remark}

\section{Adjacency Matrix}

\begin{definition}[Adjacency Matrix]
The adjacency matrix $A \in \R^{n\times n}$ is defined by
\[
a_{ij}
=
\begin{cases}
1, & \text{if vertices $i$ and $j$ are connected},\\
0, & \text{otherwise}.
\end{cases}
\]
\end{definition}

\begin{remark}
For undirected graphs, $A$ is symmetric:
\[
A = A^T.
\]
\end{remark}

\begin{example}[Adjacency Matrix]
\[
A =
\begin{pmatrix}
0 & 1 & 1 & 1 & 0\\
1 & 0 & 0 & 1 & 0\\
1 & 0 & 0 & 1 & 1\\
1 & 1 & 1 & 0 & 1\\
0 & 0 & 1 & 1 & 0
\end{pmatrix}.
\]
\end{example}

\section{Degree Matrix and Graph Laplacian}

\begin{definition}[Degree Matrix]
The degree matrix $D$ is diagonal with
\[
D_{ii} = \sum_{j=1}^n a_{ij}.
\]
\end{definition}

\begin{definition}[Graph Laplacian]
The (unnormalized) graph Laplacian is
\[
L = D - A.
\]
\end{definition}

\begin{remark}
The Laplacian encodes connectivity structure and plays a central role in spectral methods.
\end{remark}

\section{Spectral Embedding}

Spectral embedding is based on eigenvectors of the Laplacian.

\begin{definition}[Spectral Embedding Coordinates]
Let $L$ have eigenvalues
\[
0 = \lambda_1 \le \lambda_2 \le \cdots \le \lambda_n
\]
with corresponding eigenvectors $v_1,\dots,v_n$.

The spectral embedding of dimension $k$ is given by
\[
Z =
\begin{pmatrix}
v_2 & v_3 & \cdots & v_{k+1}
\end{pmatrix}.
\]
\end{definition}

\begin{remark}
The first eigenvector $v_1$ is constant and typically discarded.
\end{remark}

\section{Spectral Clustering Algorithm}

\begin{algorithm}[H]
\caption{Spectral Clustering}
\begin{algorithmic}[1]
\Require Data $\{x_i\}$, number of clusters $k$
\Ensure Cluster assignments
\State Construct similarity graph and adjacency matrix $A$
\State Compute degree matrix $D$
\State Compute Laplacian $L = D - A$
\State Compute first $k$ eigenvectors of $L$
\State Form matrix $Z$ from these eigenvectors
\State Apply $k$-means to rows of $Z$
\State \Return cluster labels
\end{algorithmic}
\end{algorithm}

\section{Interpretation}

\begin{itemize}
    \item The graph encodes local similarity structure
    \item Eigenvectors capture global structure
    \item Clustering is performed in the embedded space
\end{itemize}

\begin{remark}
Spectral methods can identify nonlinearly separable clusters that are not well captured by $k$-means in the original space.
\end{remark}

\section{Properties of the Laplacian}

\begin{proposition}
The Laplacian matrix $L$ is symmetric and positive semidefinite.
\end{proposition}

\begin{proof}
For any $x \in \R^n$,
\[
x^T L x
=
x^T(D-A)x
=
\sum_{i} d_i x_i^2 - \sum_{i,j} a_{ij} x_i x_j.
\]
Rewriting,
\[
x^T L x
=
\frac{1}{2}
\sum_{i,j} a_{ij} (x_i - x_j)^2 \ge 0.
\]
\end{proof}

\begin{proposition}
The multiplicity of the eigenvalue $0$ equals the number of connected components in the graph.
\end{proposition}

\section{Implementation Notes}

\begin{lstlisting}[language=Python]
from sklearn.cluster import SpectralClustering

model = SpectralClustering(
    n_clusters=3,
    affinity="nearest_neighbors"
)

labels = model.fit_predict(X)
\end{lstlisting}

\begin{remark}
Choice of similarity graph and its parameters (e.g., $k$ in $k$-NN or $\sigma$ in RBF kernels) strongly affects performance.
\end{remark}

\chapter{Laplacian Eigenmaps}

\section{Motivation}

Laplacian Eigenmaps is a nonlinear dimensionality reduction technique that preserves local neighborhood structure using a graph representation of the data.

\begin{remark}
Unlike PCA, which captures global variance, Laplacian Eigenmaps focuses on preserving local geometry.
\end{remark}

\section{Graph Construction}

Given data $\{x_i\}_{i=1}^n \subset \R^d$, construct a similarity graph:
\begin{itemize}
    \item Connect points using $k$-nearest neighbors or $\varepsilon$-neighborhoods
    \item Assign weights $w_{ij}$ to edges
\end{itemize}

A common choice is the heat kernel:
\[
w_{ij}
=
\exp\left(-\frac{\norm{x_i - x_j}^2}{2\sigma^2}\right).
\]

\begin{definition}[Weight Matrix]
Let $W \in \R^{n \times n}$ be defined by
\[
W_{ij} = w_{ij}.
\]
\end{definition}

\section{Graph Laplacian}

\begin{definition}[Degree Matrix]
\[
D_{ii} = \sum_{j=1}^n W_{ij}.
\]
\end{definition}

\begin{definition}[Laplacian Matrix]
\[
L = D - W.
\]
\end{definition}

\section{Optimization Problem}

We seek a low-dimensional embedding $y_i \in \R^m$ for each data point $x_i$.

\begin{definition}[Laplacian Eigenmaps Objective]
\[
\min_{\{y_i\}}
\quad
\frac{1}{2}
\sum_{i,j} W_{ij} \norm{y_i - y_j}_2^2.
\]
\end{definition}

\begin{remark}
Nearby points (large $W_{ij}$) are encouraged to remain close in the embedding.
\end{remark}

\subsection{Matrix Form}

Let $Y = (y_1,\dots,y_n)^T \in \R^{n \times m}$. Then
\[
\sum_{i,j} W_{ij} \norm{y_i - y_j}_2^2
=
2 \, \tr(Y^T L Y).
\]

\begin{proof}
Expand:
\[
\norm{y_i - y_j}_2^2
=
y_i^T y_i - 2 y_i^T y_j + y_j^T y_j.
\]
Summing,
\[
\sum_{i,j} W_{ij} \norm{y_i - y_j}_2^2
=
2 \sum_i d_i y_i^T y_i - 2 \sum_{i,j} W_{ij} y_i^T y_j.
\]
This equals
\[
2 \, \tr(Y^T D Y) - 2 \, \tr(Y^T W Y)
=
2 \, \tr(Y^T L Y).
\]
\end{proof}

\section{Constraints}

To avoid trivial solutions (e.g., $y_i=0$), impose constraints:
\[
Y^T D Y = I,
\quad
Y^T D \mathbf{1} = 0.
\]

\begin{remark}
The second constraint removes the constant eigenvector.
\end{remark}

\section{Solution via Eigenvalue Problem}

The optimization reduces to solving the generalized eigenvalue problem:
\[
L v = \lambda D v.
\]

\begin{definition}[Embedding Coordinates]
The embedding is given by the eigenvectors corresponding to the smallest nonzero eigenvalues:
\[
Y = [v_2,\dots,v_{m+1}].
\]
\end{definition}

\begin{remark}
The first eigenvector $v_1$ corresponds to eigenvalue $0$ and is constant.
\end{remark}

\section{Interpretation}

\begin{itemize}
    \item Minimizes distortion of local neighborhoods
    \item Preserves manifold structure
    \item Produces a low-dimensional embedding reflecting intrinsic geometry
\end{itemize}

\section{Connection to Spectral Clustering}

\begin{remark}
Laplacian Eigenmaps provides the embedding step used in spectral clustering. After computing $Y$, clustering methods (e.g., $k$-means) can be applied to the embedded coordinates.
\end{remark}

\section{Algorithm}

\begin{algorithm}[H]
\caption{Laplacian Eigenmaps}
\begin{algorithmic}[1]
\Require Data $\{x_i\}$, embedding dimension $m$
\Ensure Embedded coordinates $\{y_i\}$
\State Construct similarity graph and weight matrix $W$
\State Compute degree matrix $D$
\State Compute Laplacian $L = D - W$
\State Solve $L v = \lambda D v$
\State Take eigenvectors $v_2,\dots,v_{m+1}$
\State Form embedding $Y = [v_2,\dots,v_{m+1}]$
\State \Return $\{y_i\}$
\end{algorithmic}
\end{algorithm}

\section{Properties and Limitations}

\begin{itemize}
    \item Captures nonlinear structure
    \item Sensitive to graph construction (choice of $k$, $\sigma$)
    \item Requires eigen-decomposition (computational cost)
\end{itemize}

\section{Implementation Notes}

\begin{lstlisting}[language=Python]
from sklearn.manifold import SpectralEmbedding

model = SpectralEmbedding(
    n_components=2,
    n_neighbors=10
)

Y = model.fit_transform(X)
\end{lstlisting}

\begin{remark}
The embedding can be visualized or used as input to clustering or classification algorithms.
\end{remark}

\chapter{Manifold Learning}

\section{Motivation}

Manifold learning assumes that high-dimensional data lie on or near a low-dimensional manifold embedded in $\R^D$.

\begin{remark}
The goal is to recover intrinsic low-dimensional coordinates while preserving geometric structure (distances, neighborhoods, or local linearity).
\end{remark}

\section{Metric Multidimensional Scaling (MDS)}

\subsection{Problem Formulation}

Given a distance matrix $\{d_{ij}\}$, find points $\{x_i\}_{i=1}^n \subset \R^d$ such that
\[
d_{ij} \approx \norm{x_i - x_j}_2.
\]

\begin{definition}[MDS Stress Function]
\[
\min_{\{x_i\}}
\quad
\sum_{i<j}
\left(
\norm{x_i - x_j}_2 - d_{ij}
\right)^2.
\]
\end{definition}

\subsection{Classical MDS}

Define the squared distance matrix $D^{(2)}$ with entries $d_{ij}^2$. Let
\[
J = I - \frac{1}{n}\mathbf{1}\mathbf{1}^T.
\]

\begin{definition}[Centered Gram Matrix]
\[
B = -\frac{1}{2} J D^{(2)} J.
\]
\end{definition}

\begin{proposition}
If $B = X X^T$, then the rows of $X$ provide the embedding coordinates.
\end{proposition}

\begin{remark}
The embedding is obtained via eigen-decomposition:
\[
B = U \Lambda U^T,
\quad
X = U_d \Lambda_d^{1/2}.
\]
\end{remark}

\section{Isomap}

\subsection{Idea}

Isomap extends MDS by replacing Euclidean distances with geodesic distances on a neighborhood graph.

\begin{remark}
Geodesic distances approximate distances along the manifold rather than through ambient space.
\end{remark}

\subsection{Algorithm}

\begin{algorithm}[H]
\caption{Isomap}
\begin{algorithmic}[1]
\Require Data $\{x_i\}$, neighborhood size ($k$ or $\varepsilon$), embedding dimension $d$
\Ensure Embedded coordinates
\State Construct neighborhood graph:
\begin{itemize}
    \item Connect $k$ nearest neighbors or points within $\varepsilon$
\end{itemize}
\State Initialize graph distances:
\[
d_G(i,j) =
\begin{cases}
\norm{x_i - x_j}_2, & \text{if connected},\\
\infty, & \text{otherwise}.
\end{cases}
\]
\State Compute shortest paths:
\[
d_G(i,j) = \min_k \left\{ d_G(i,j), \; d_G(i,k) + d_G(k,j) \right\}
\]
\State Apply classical MDS to $\{d_G(i,j)\}$
\State \Return embedding
\end{algorithmic}
\end{algorithm}

\begin{remark}
Shortest paths can be computed using Floyd–Warshall or Dijkstra's algorithm.
\end{remark}

\section{Locally Linear Embedding (LLE)}

\subsection{Idea}

LLE assumes each point can be reconstructed as a linear combination of its neighbors.

\begin{definition}[Reconstruction Weights]
For each point $x_i$, find weights $w_{ij}$ solving
\[
\min_{w_{ij}}
\quad
\norm{x_i - \sum_{j \in \mathcal{N}(i)} w_{ij} x_j}_2^2
\]
subject to
\[
\sum_{j \in \mathcal{N}(i)} w_{ij} = 1.
\]
\end{definition}

\subsection{Embedding Step}

Find low-dimensional points $\{y_i\}$ minimizing
\[
\sum_i
\norm{y_i - \sum_j w_{ij} y_j}_2^2.
\]

\begin{remark}
The weights $w_{ij}$ are fixed from the original space and reused in the embedding space.
\end{remark}

\subsection{Solution}

This reduces to an eigenvalue problem:
\[
M = (I - W)^T (I - W),
\]
and the embedding is given by the eigenvectors corresponding to the smallest nonzero eigenvalues.

\section{Comparison of Methods}

\begin{itemize}
    \item MDS:
    \begin{itemize}
        \item Preserves pairwise distances
        \item Sensitive to noise in distances
    \end{itemize}
    \item Isomap:
    \begin{itemize}
        \item Preserves geodesic distances
        \item Captures global manifold structure
    \end{itemize}
    \item LLE:
    \begin{itemize}
        \item Preserves local linear structure
        \item Does not explicitly preserve distances
    \end{itemize}
\end{itemize}

\section{Other Embedding Methods}

\subsection{Random Projection}

Project data using a random matrix:
\[
y = R x,
\]
where $R$ has random entries.

\begin{remark}
Preserves distances approximately (Johnson–Lindenstrauss lemma).
\end{remark}

\subsection{PCA Projection}

Project onto top principal components:
\[
y = E_K^T x.
\]

\subsection{Spectral Embedding}

Uses eigenvectors of a graph Laplacian to embed data.

\subsection{$t$-SNE}

\begin{remark}
$t$-SNE maps high-dimensional data to low dimensions by preserving local similarities using probability distributions.
\end{remark}

\section{Implementation Notes}

\begin{lstlisting}[language=Python]
from sklearn.manifold import MDS, Isomap, LocallyLinearEmbedding

mds = MDS(n_components=2)
X_mds = mds.fit_transform(X)

iso = Isomap(n_neighbors=10, n_components=2)
X_iso = iso.fit_transform(X)

lle = LocallyLinearEmbedding(n_neighbors=10, n_components=2)
X_lle = lle.fit_transform(X)
\end{lstlisting}

\begin{lstlisting}[language=Python]
from sklearn.manifold import TSNE

tsne = TSNE(n_components=2)
X_tsne = tsne.fit_transform(X)
\end{lstlisting}

\begin{remark}
Choice of method depends on whether global or local structure is more important.
\end{remark}

\chapter{SQLite Databases and SQL}

\section{Overview}

Relational databases store data in structured tables and are queried using Structured Query Language (SQL).

\begin{remark}
SQLite is a lightweight, self-contained database engine that can be embedded directly in applications (e.g., Python).
\end{remark}

\section{Database Schema}

We consider a synthetic measurement database with four tables:

\begin{itemize}
    \item \textbf{Data(X, Y, RunID, ID)}: measurement points
    \item \textbf{Runs(RunID, UserID, InstrumentID)}: run metadata
    \item \textbf{Users(UserID, Name)}: user information
    \item \textbf{Instruments(InsID, Name)}: instrument information
\end{itemize}

\begin{remark}
Tables are related via keys:
\begin{itemize}
    \item RunID links Data and Runs
    \item UserID links Runs and Users
    \item InstrumentID links Runs and Instruments
\end{itemize}
\end{remark}

\section{Basic SQL Queries}

\subsection{Selecting Data}

\begin{lstlisting}[language=SQL]
SELECT * FROM Data;
\end{lstlisting}

\begin{lstlisting}[language=SQL]
SELECT X, Y FROM Data WHERE X > 0;
\end{lstlisting}

\subsection{Sorting}

\begin{lstlisting}[language=SQL]
SELECT * FROM Data
ORDER BY X DESC;
\end{lstlisting}

\section{Aggregation}

\begin{definition}[Aggregate Functions]
SQL provides functions that summarize data:
\begin{itemize}
    \item COUNT
    \item SUM
    \item AVG
    \item MIN
    \item MAX
\end{itemize}
\end{definition}

\begin{example}
\begin{lstlisting}[language=SQL]
SELECT COUNT(*) FROM Data;

SELECT AVG(X), AVG(Y) FROM Data;
\end{lstlisting}
\end{example}

\subsection{Grouping}

\begin{lstlisting}[language=SQL]
SELECT RunID, COUNT(*) AS n_points
FROM Data
GROUP BY RunID;
\end{lstlisting}

\begin{remark}
GROUP BY partitions rows and applies aggregate functions within each group.
\end{remark}

\section{Joins}

\begin{definition}[Join]
A join combines rows from multiple tables based on matching keys.
\end{definition}

\subsection{Inner Join}

\begin{lstlisting}[language=SQL]
SELECT *
FROM Data d
JOIN Runs r ON d.RunID = r.RunID;
\end{lstlisting}

\subsection{Multiple Joins}

\begin{lstlisting}[language=SQL]
SELECT u.Name, i.Name, COUNT(*) AS n
FROM Data d
JOIN Runs r ON d.RunID = r.RunID
JOIN Users u ON r.UserID = u.UserID
JOIN Instruments i ON r.InstrumentID = i.InsID
GROUP BY u.Name, i.Name;
\end{lstlisting}

\begin{remark}
Joins enable combining relational data across tables.
\end{remark}

\section{Subqueries}

\begin{definition}[Subquery]
A subquery is a query nested inside another query.
\end{definition}

\begin{example}
\begin{lstlisting}[language=SQL]
SELECT *
FROM Data
WHERE RunID IN (
    SELECT RunID FROM Runs WHERE UserID = 1
);
\end{lstlisting}
\end{example}

\section{Common Table Expressions (CTEs)}

\begin{definition}[CTE]
A Common Table Expression defines a temporary result set for use within a query.
\end{definition}

\begin{example}
\begin{lstlisting}[language=SQL]
WITH RunCounts AS (
    SELECT RunID, COUNT(*) AS n
    FROM Data
    GROUP BY RunID
)
SELECT *
FROM RunCounts
WHERE n > 10;
\end{lstlisting}
\end{example}

\begin{remark}
CTEs improve readability and modularity of complex queries.
\end{remark}

\section{Relational Algebra Perspective}

SQL operations correspond to relational algebra:

\begin{itemize}
    \item SELECT (projection): choose columns
    \item WHERE (selection): filter rows
    \item JOIN: combine relations
    \item GROUP BY: aggregation
\end{itemize}

\begin{remark}
Relational algebra provides a formal foundation for database queries.
\end{remark}

\section{Table Creation and Modification}

\subsection{Creating Tables}

\begin{lstlisting}[language=SQL]
CREATE TABLE Users (
    UserID INTEGER PRIMARY KEY,
    Name TEXT
);
\end{lstlisting}

\subsection{Inserting Data}

\begin{lstlisting}[language=SQL]
INSERT INTO Users (UserID, Name)
VALUES (1, 'Alice');
\end{lstlisting}

\subsection{Updating Data}

\begin{lstlisting}[language=SQL]
UPDATE Users
SET Name = 'Bob'
WHERE UserID = 1;
\end{lstlisting}

\subsection{Deleting Data}

\begin{lstlisting}[language=SQL]
DELETE FROM Users
WHERE UserID = 1;
\end{lstlisting}

\section{Transactions}

\begin{definition}[Transaction]
A transaction is a sequence of operations executed as a single unit.
\end{definition}

\begin{lstlisting}[language=SQL]
BEGIN TRANSACTION;

INSERT INTO Data (X, Y, RunID)
VALUES (1.0, 2.0, 3);

COMMIT;
\end{lstlisting}

\begin{remark}
Transactions ensure:
\begin{itemize}
    \item Atomicity
    \item Consistency
    \item Isolation
    \item Durability (ACID properties)
\end{itemize}
\end{remark}

\section{Python Integration}

\begin{lstlisting}[language=Python]
import sqlite3

conn = sqlite3.connect(":memory:")
cursor = conn.cursor()

cursor.execute("SELECT * FROM Data")
rows = cursor.fetchall()

conn.close()
\end{lstlisting}

\begin{remark}
SQLite integrates seamlessly with Python for building lightweight data pipelines and experiments.
\end{remark}

\end{document}